\documentclass[a4paper,12pt,twoside]{memoir}

% Castellano
\usepackage[spanish,es-tabla]{babel}
\selectlanguage{spanish}
\usepackage[utf8]{inputenc}
\usepackage[T1]{fontenc}
\usepackage{lmodern} % Scalable font
\usepackage{microtype}
\usepackage{placeins}

\RequirePackage{booktabs}
\RequirePackage[table]{xcolor}
\RequirePackage{xtab}
\RequirePackage{multirow}

% Links
\PassOptionsToPackage{hyphens}{url}\usepackage[colorlinks]{hyperref}
\hypersetup{
	allcolors = {red}
}

% Ecuaciones
\usepackage{amsmath}

% Rutas de fichero / paquete
\newcommand{\ruta}[1]{{\sffamily #1}}

% Párrafos
\nonzeroparskip

% Huérfanas y viudas
\widowpenalty100000
\clubpenalty100000

% Imagenes
\usepackage{graphicx}
\newcommand{\imagen}[2]{
	\begin{figure}[!h]
		\centering
		\includegraphics[width=0.9\textwidth]{#1}
		\caption{#2}\label{fig:#1}
	\end{figure}
	\FloatBarrier
}

\newcommand{\imagenflotante}[2]{
	\begin{figure}%[!h]
		\centering
		\includegraphics[width=0.9\textwidth]{#1}
		\caption{#2}\label{fig:#1}
	\end{figure}
}



% El comando \figura nos permite insertar figuras comodamente, y utilizando
% siempre el mismo formato. Los parametros son:
% 1 -> Porcentaje del ancho de página que ocupará la figura (de 0 a 1)
% 2 --> Fichero de la imagen
% 3 --> Texto a pie de imagen
% 4 --> Etiqueta (label) para referencias
% 5 --> Opciones que queramos pasarle al \includegraphics
% 6 --> Opciones de posicionamiento a pasarle a \begin{figure}
\newcommand{\figuraConPosicion}[6]{%
  \setlength{\anchoFloat}{#1\textwidth}%
  \addtolength{\anchoFloat}{-4\fboxsep}%
  \setlength{\anchoFigura}{\anchoFloat}%
  \begin{figure}[#6]
    \begin{center}%
      \Ovalbox{%
        \begin{minipage}{\anchoFloat}%
          \begin{center}%
            \includegraphics[width=\anchoFigura,#5]{#2}%
            \caption{#3}%
            \label{#4}%
          \end{center}%
        \end{minipage}
      }%
    \end{center}%
  \end{figure}%
}

%
% Comando para incluir imágenes en formato apaisado (sin marco).
\newcommand{\figuraApaisadaSinMarco}[5]{%
  \begin{figure}%
    \begin{center}%
    \includegraphics[angle=90,height=#1\textheight,#5]{#2}%
    \caption{#3}%
    \label{#4}%
    \end{center}%
  \end{figure}%
}
% Para las tablas
\newcommand{\otoprule}{\midrule [\heavyrulewidth]}
%
% Nuevo comando para tablas pequeñas (menos de una página).
\newcommand{\tablaSmall}[5]{%
 \begin{table}
  \begin{center}
   \rowcolors {2}{gray!35}{}
   \begin{tabular}{#2}
    \toprule
    #4
    \otoprule
    #5
    \bottomrule
   \end{tabular}
   \caption{#1}
   \label{tabla:#3}
  \end{center}
 \end{table}
}

%
% Nuevo comando para tablas pequeñas (menos de una página).
\newcommand{\tablaSmallSinColores}[5]{%
 \begin{table}[H]
  \begin{center}
   \begin{tabular}{#2}
    \toprule
    #4
    \otoprule
    #5
    \bottomrule
   \end{tabular}
   \caption{#1}
   \label{tabla:#3}
  \end{center}
 \end{table}
}

\newcommand{\tablaApaisadaSmall}[5]{%
\begin{landscape}
  \begin{table}
   \begin{center}
    \rowcolors {2}{gray!35}{}
    \begin{tabular}{#2}
     \toprule
     #4
     \otoprule
     #5
     \bottomrule
    \end{tabular}
    \caption{#1}
    \label{tabla:#3}
   \end{center}
  \end{table}
\end{landscape}
}

%
% Nuevo comando para tablas grandes con cabecera y filas alternas coloreadas en gris.
\newcommand{\tabla}[6]{%
  \begin{center}
    \tablefirsthead{
      \toprule
      #5
      \otoprule
    }
    \tablehead{
      \multicolumn{#3}{l}{\small\sl continúa desde la página anterior}\\
      \toprule
      #5
      \otoprule
    }
    \tabletail{
      \hline
      \multicolumn{#3}{r}{\small\sl continúa en la página siguiente}\\
    }
    \tablelasttail{
      \hline
    }
    \bottomcaption{#1}
    \rowcolors {2}{gray!35}{}
    \begin{xtabular}{#2}
      #6
      \bottomrule
    \end{xtabular}
    \label{tabla:#4}
  \end{center}
}

%
% Nuevo comando para tablas grandes con cabecera.
\newcommand{\tablaSinColores}[6]{%
  \begin{center}
    \tablefirsthead{
      \toprule
      #5
      \otoprule
    }
    \tablehead{
      \multicolumn{#3}{l}{\small\sl continúa desde la página anterior}\\
      \toprule
      #5
      \otoprule
    }
    \tabletail{
      \hline
      \multicolumn{#3}{r}{\small\sl continúa en la página siguiente}\\
    }
    \tablelasttail{
      \hline
    }
    \bottomcaption{#1}
    \begin{xtabular}{#2}
      #6
      \bottomrule
    \end{xtabular}
    \label{tabla:#4}
  \end{center}
}

%
% Nuevo comando para tablas grandes sin cabecera.
\newcommand{\tablaSinCabecera}[5]{%
  \begin{center}
    \tablefirsthead{
      \toprule
    }
    \tablehead{
      \multicolumn{#3}{l}{\small\sl continúa desde la página anterior}\\
      \hline
    }
    \tabletail{
      \hline
      \multicolumn{#3}{r}{\small\sl continúa en la página siguiente}\\
    }
    \tablelasttail{
      \hline
    }
    \bottomcaption{#1}
  \begin{xtabular}{#2}
    #5
   \bottomrule
  \end{xtabular}
  \label{tabla:#4}
  \end{center}
}



\definecolor{cgoLight}{HTML}{EEEEEE}
\definecolor{cgoExtralight}{HTML}{FFFFFF}

%
% Nuevo comando para tablas grandes sin cabecera.
\newcommand{\tablaSinCabeceraConBandas}[5]{%
  \begin{center}
    \tablefirsthead{
      \toprule
    }
    \tablehead{
      \multicolumn{#3}{l}{\small\sl continúa desde la página anterior}\\
      \hline
    }
    \tabletail{
      \hline
      \multicolumn{#3}{r}{\small\sl continúa en la página siguiente}\\
    }
    \tablelasttail{
      \hline
    }
    \bottomcaption{#1}
    \rowcolors[]{1}{cgoExtralight}{cgoLight}

  \begin{xtabular}{#2}
    #5
   \bottomrule
  \end{xtabular}
  \label{tabla:#4}
  \end{center}
}



\graphicspath{ {./img/} }

% Capítulos
\chapterstyle{bianchi}
\newcommand{\capitulo}[2]{
	\setcounter{chapter}{#1}
	\setcounter{section}{0}
	\setcounter{figure}{0}
	\setcounter{table}{0}
	\chapter*{#2}
	\addcontentsline{toc}{chapter}{#2}
	\markboth{#2}{#2}
}

% Apéndices
\renewcommand{\appendixname}{Apéndice}
\renewcommand*\cftappendixname{\appendixname}

\newcommand{\apendice}[1]{
	%\renewcommand{\thechapter}{A}
	\chapter{#1}
}

\renewcommand*\cftappendixname{\appendixname\ }

% Formato de portada
\makeatletter
\usepackage{xcolor}
\newcommand{\tutor}[1]{\def\@tutor{#1}}
\newcommand{\course}[1]{\def\@course{#1}}
\definecolor{cpardoBox}{HTML}{E6E6FF}
\def\maketitle{
  \null
  \thispagestyle{empty}
  % Cabecera ----------------
\noindent\includegraphics[width=\textwidth]{cabecera}\vspace{1cm}%
  \vfill
  % Título proyecto y escudo informática ----------------
  \colorbox{cpardoBox}{%
    \begin{minipage}{.8\textwidth}
      \vspace{.5cm}\Large
      \begin{center}
      \textbf{TFG del Grado en Ingeniería Informática}\vspace{.6cm}\\
      \textbf{\LARGE\@title{}}
      \end{center}
      \vspace{.2cm}
    \end{minipage}

  }%
  \hfill\begin{minipage}{.20\textwidth}
    \includegraphics[width=\textwidth]{escudoInfor}
  \end{minipage}
  \vfill
  % Datos de alumno, curso y tutores ------------------
  \begin{center}%
  {%
    \noindent\LARGE
    Presentado por \@author{}\\ 
    en Universidad de Burgos --- \@date{}\\
    Tutor: \@tutor{}\\
  }%
  \end{center}%
  \null
  \cleardoublepage
  }
\makeatother

\newcommand{\nombre}{Nombre del alumno} %%% cambio de comando

% Datos de portada
\title{título del TFG}
\author{\nombre}
\tutor{nombre tutor}
\date{\today}

\begin{document}

\maketitle


\newpage\null\thispagestyle{empty}\newpage


%%%%%%%%%%%%%%%%%%%%%%%%%%%%%%%%%%%%%%%%%%%%%%%%%%%%%%%%%%%%%%%%%%%%%%%%%%%%%%%%%%%%%%%%
\thispagestyle{empty}


\noindent\includegraphics[width=\textwidth]{cabecera}\vspace{1cm}

\noindent D. nombre tutor, profesor del departamento de nombre departamento, área de nombre área.

\noindent Expone:

\noindent Que el alumno D. \nombre, con DNI dni, ha realizado el Trabajo final de Grado en Ingeniería Informática titulado título de TFG. 

\noindent Y que dicho trabajo ha sido realizado por el alumno bajo la dirección del que suscribe, en virtud de lo cual se autoriza su presentación y defensa.

\begin{center} %\large
En Burgos, {\large \today}
\end{center}

\vfill\vfill\vfill

% Author and supervisor
\begin{minipage}{0.45\textwidth}
\begin{flushleft} %\large
Vº. Bº. del Tutor:\\[2cm]
D. nombre tutor
\end{flushleft}
\end{minipage}
\hfill
\begin{minipage}{0.45\textwidth}
\begin{flushleft} %\large
Vº. Bº. del co-tutor:\\[2cm]
D. nombre co-tutor
\end{flushleft}
\end{minipage}
\hfill

\vfill

% para casos con solo un tutor comentar lo anterior
% y descomentar lo siguiente
%Vº. Bº. del Tutor:\\[2cm]
%D. nombre tutor


\newpage\null\thispagestyle{empty}\newpage




\frontmatter

% Abstract en castellano
\renewcommand*\abstractname{Resumen}
\begin{abstract}
En este primer apartado se hace una \textbf{breve} presentación del tema que se aborda en el proyecto.
\end{abstract}

\renewcommand*\abstractname{Descriptores}
\begin{abstract}
Palabras separadas por comas que identifiquen el contenido del proyecto Ej: servidor web, buscador de vuelos, android \ldots
\end{abstract}

\clearpage

% Abstract en inglés
\renewcommand*\abstractname{Abstract}
\begin{abstract}
A \textbf{brief} presentation of the topic addressed in the project.
\end{abstract}

\renewcommand*\abstractname{Keywords}
\begin{abstract}
keywords separated by commas.
\end{abstract}

\clearpage

% Indices
\tableofcontents

\clearpage

\listoffigures

\clearpage

\listoftables
\clearpage

\mainmatter
\capitulo{1}{Introducción}

En este proyecto vamos a implementar una transcrpción automática de lenguaje de signos a texto a nivel palabra usando técnicas de inteligencia artificial.

El lenguaje de signos es una forma de comunicación por gestos gracias a la cual más de 72 millones de personas con disacapacidad auditiva pueden comunicarse. Hay más de 300 lenguas de signo distintas y dentro de un mismo lenguaje hay varios gestos que pueden significar palabras distintas dependiendo del contexto.

En la actualidad existen transcriptores a nivel letra o nivel palabra, solo que estos transcriptores solo detectan gestos estáticos, que no tienen movimiento. No existe ningún traductor de lenguaje de signos que sea capaz de transcribir gestos con movimiento incluido.

\project{Sign2Text} tiene el propósito de crear un transcriptor a tiempo real usando redes neuronales, \loc{deep learning} y procesado de imagen y video. Para esto se ha usado un dataset muy completo de \sigla{ASL} (\loc{American Sign Language}) \bib{li2020word} compuesto por secuencias cortas de video e información de poses de distintos sujetos. Se van a aplicar técnicas de \loc{deep learning} como redes neuronales convolucionales, \loc{transformers} y procesado 3D (dos dimensiones para la imagen y otra para el tiempo).

Tras crear la red que transcriba el texto, se pretende crear una demo \loc{web} para que cualquier usuario alrededor del mundo pueba probarlo. También se busca poder contenedorizar el modelo y dejarlo libre para que cualquier desarrollador pueda descargarlo y crear su propia aplicación.


\section{Estructura del trabajo}

\begin{itemize}
  \item \textbf{Objetivos del trabajo}: En este apartado vamos a listar los objetivos que queremos cumplir en la realización del proyecto, así como las caracterísitcas que esperamos cumpla nuestro modelo de \loc{deep learning}.
  \item \textbf{Conceptos teóricos}: En este apartado vamos a hacer un \loc{overview} de \loc{machine learning}, comenzando por el aprendizaje no supervisado, siguiendo del semi-supervisado y terminando en lo que más no interesa, el supervisado y el \loc{deep learning}. Veremos las prinicipales técnicas en cada tipo de aprendizaje para conocer cuales son las mejores opciones en el desarrollo del proyecto.
  \item \textbf{Técnicas y herramientas}: Aquí listaremos y explicaremos todas las herramientas usadas en la realización de \project{Sign2Text}. También comentaremos un poco como las hemos usado nosotros y que alternativas tenemos en cada caso. Hablaremos de las herramientas con las que desarrollaremos nuestro modelo y con las que crearemos nuestro \loc{front-end} y \loc{back-end}.
  \item \textbf{Aspectos relevantes del desarrollo}: En este apartado hablaremos sobre exprimentos y experiencias del proyecto. Aquí contaremos las mejoras que se han hecho a media que iteramos con el modelo, así como curiosidades y problemas con los que nos hayamos encontrado.
  \item \textbf{Trabajos relacionados}: Esta sección pretende mostrar otras herramientas y proyectos similares a \project{Sign2Text} disponibles en la comunidad y el mercado.
  \item \textbf{Conclusiones y líneas futuras}: Por último haremos un repaso de las mejoras que se pueden realizar en el proyecto. Hablaremos de las conclusiones que hemos sacado al finalizarlo y repasaremos que porcentaje de los objetivos hemos cumplido al finalizar el proyecto.
\end{itemize}
\capitulo{2}{Objetivos del proyecto}

La principal motivación del proyecto es:

\textbf{
  Mejorar la interacción entre personas que necesitan comunicación por signos y personas que no conocen el lenguaje de signos, intentando aumentar la calidad de vida de las primeras.
}


\section{Objetivos teóricos}

\begin{enumerate}
  \item Crear de un modelo de \sigla{DL}\footnote{\textit{Deep Learning}} capaz de clasificar ASL\footnote{\textit{American Sign Language}} a nivel palabra.
  \item Poner en producción un método para que los usuarios puedan probar el modelo de forma libre y transparente.
  \item Generar una \sigla{API} de código abierto para que otros desarrolladores puedan crear plataformas de cliente, con el objeto de aumentar la audiencia y alcance del proyecto.
  \item Hacer disponible y encapsular el modelo en un contenedor Docker para su libre distribución.
  \item Crear de herramientas que permitan el tratamiento de datos y la transformación entre distintos formatos\footnote{\pe: extracción de fotogramas de un vídeo, concatenación de imágenes}.
  \item Estudiar del comportamiento, estructura e hiperparametrización de redes neuronales convolucionales.
  \item Aprender a fondo el uso de \prog{PyTorch}\cite{PYTORCH}, un \loc{framework} para acelerar la creación y prototipado de redes neuronales.
  \item Exportar del modelo entrenado a un formato estándar que facilite la compatibilidad con el mayor número de librerías en distintos lenguajes de programación.
  \item Mantener de una \loc{codebase} que favorezca la integración continua (\loc{linting}\footnote{Voz ingl. Estudio de limpieza, orden, calidad y redundancia}, \loc{formatting}\footnote{Formatear: Correcta estructura y legibilidad} y \loc{type-checking}\footnote{Comprobación de tipados de variables y funciones}) de código, siguiendo así los estándares en contribuciones \loc{Open Source}.
\end{enumerate}


Estos objetivos representan en general la filosofía y los objetivos teóricos del proyecto. Si entramos más en detalle sobre los aspectos técnicos y las \loc{features} que se esperan obtener al finalizar el proyecto, obtenemos los siguientes puntos:

\section{Caracterísitcas esperadas}

\begin{itemize}
  \item el sistema debe ser capaz de inferir resultados del modelo entrenado a tiempo real.
  \item Se desarrollará un modelo fácilmente escalable y adaptable a distintos dataset de cualquier tamaño.\footnote{El formato de los \loc{datasets} debe ser el mismo. Un modelo entrenado para clasificación de imágenes no podrá clasificar en formato vídeo o audio.}
  \item Se realizará una fase de \loc{data augmentation} sobre los datos iniciales.
  \item Se implementarán distintas redes neuronales para los distintos formatos de datos (imagen y vídeo). Deben ser fácilmente refactorizables y mantenibles.
  \item Cada vez que se estudie el uso de un nuevo formato de dataset, se creará una nueva red, manteniendo la usabilidad de la anterior intacta.
  \item A lo largo de las pruebas y entrenamientos de la red, vamos a probar distintos \loc{schedulers}, \loc{optimizers} y \loc{criterions}\footnote{Funciones de cálculo de pérdida} buscando el que mejor se adecúe al formato de los datos y la estructura de la red.
  \item Se mantendrán unas estadísticas a tiempo real de la fase de \loc{training} y \loc{test} mediante el uso de \prog{Tensorboard}, una herramienta analítica que permite mantener un \loc{log} de imágenes generadas en la ejecución; así como gráficas de costes y \loc{accuracies}.
\end{itemize}

\capitulo{3}{Conceptos teóricos}

%%%%%%%%%%%%%%%%%%%%%%%%%%%%%%%%%%%%%%%%%%%%%%%%%%%%%%%%%%%%%

Este proyecto es un proyecto de \loc{machine learning}. Dentro del \loc{machine learning}, es de aprendizaje supervisado. Más concretamente intenta resolver un problema de clasificación mediante \loc{deep learning}. Las redes neuronales aplicadas para la clasificación serán \sigla{CNN} (\loc{Convolutional Neural Networks}), y \loc{ResNets}, una arquitectura sobre las \sigla{CNN}. Veamos en detalle todos estos términos en los siguientes apartados.

\section{Machine Learning}

Según Tom Mitchell \bib{mitchell1997machine}, un programa de ordenador aprende de una experiencia \textit{E} con respecto a alguna tarea \textit{T} si su medida de desempeño \textit{P} (del inglés \loc{performance}) mejora con la experiencia E.

El uso del machine learning se origina en la búsqueda de soluciones a problemas que no podemos resolver programáticamente. Determinar las soluciones se centra en la recogida y análisis de datos\footnote{A lo largo del documento nos referiremos a los datos indistintamente como <<datos>> o <<instancias>>} con la finalidad de buscar relaciones entre ellos.

Para aclarar este concepto, veamos un ejemplo \bib{alpaydin2021machine}\label{example:car_price}:
Imaginemos que debemos estimar el precio de un coche usado pero no tenemos la fórmula exacta para valorar su estado. Lo que si sabemos es que el precio del coche aumenta y disminuye dependiendo de sus propiedades, como la marca, el kilometraje, o el desgaste. No sabemos tampoco cuales de las propiedades tienen más importancia. Sabemos seguro que cuantos más kilómetros tenga el coche, menor será su precio, pero no sabemos a que escala ocurre esto.

Para determinar la solución necesitamos estudiar el precio de coches en el mercado y anotar sus características. Estos datos son los que daremos a nuestro programa para que busque las relaciones entre propiedades e indique el precio óptimo de los coches.

Saliendo del ejemplo, el \loc{machine learning} se divide en distintos tipos según la cantidad de datos (o supervisión) que otorgamos al <<programa>>: \loc{supervised, unsupervised} y \loc{semi-supervised learning}. Hay más formas de dividir el \loc{machine learning}, como puedes ver en \ref{fig:MachineLearningTypes} (nombraremos alguna de ella a lo largo de los siguientes apartados), pero nos vamos a centrar únicamente en \loc{supervised learning}.

\imagen{./img/memoria/conceptos/MachineLearningTypes}{Tipos de \loc{machine learning} según datos, inferencia estadística y técnicas de aprendizaje}{MachineLearningTypes}

% \subsection{Unsupervised Learning}

% El \loc{unsupervised learning} o  aprendizaje no supervisado, tiene su principal característica en que es capaz de detectar las relaciones entre los datos sin ayuda de etiquetas o respuestas externas. Si volvemos al ejemplo de los coches, seríamos capaces de determinar el precio de venta sin necesidad de recibir datos de coches reales.

% La técnica principal en el aprendizaje no supervisado es el agrupamiento de los datos (o \loc{clustering}) según la similaridad de sus características.

% \subsubsection{\loc{Clustering}}

% El \loc{clustering} consiste en el agrupamiento de datos no etiquetados basándonos en sus similaridades o diferencias. Los algoritmos de \loc{clustering} son usados para procesar datos crudos sin alteraciones en grupos representados por patrones o estructuras de información. Estos algoritmos se pueden clasificar en:

% \begin{itemize}
%   \item \textbf{Mutuamente excluyentes}: En estos algoritmos se estipula que una instancia de datos solo puede existir en un y solo un \loc{cluster} (o grupo). Este tipo de \loc{clustering} también se denomina \loc{<<hard>> clustering}. Un ejemplo de algoritmo mutuamente excluyente es el algoritmo de \loc{K-means}. En \loc{K-means}, la <<K>> indica el número de \loc{clusters}  basándonos en su centroide\footnote{O centro geométrico, es la posición media aritmética entre todos los puntos de una figura}. Los puntos, representando a un dato, caerán en el grupo con el centroide más cercano. Un valor de \loc{K-means} mayor indicará un mayor número de agrupamientos. \imagen{./img/memoria/conceptos/Unsupervised_K-Means_Clustering_Algorithm}{Ejemplo de datos antes y después del \loc{clustering} para un algoritmo mutuamente excluyente \loc{K-means} \bib{sinaga2020unsupervised}.}{BeforeAfterClustering}

%   \item \textbf{Superpuestos}: En el caso de los algoritmos superpuestos, la diferencia es que un dato puede pertenecer a más de un \loc{cluster} de forma simultánea con distinto grado de pertenencia. Un ejemplo de algoritmo superpuesto es el \loc{<<soft>> K-means}.

%   \item \textbf{Jerárquicos}: También conocido como \sigla{HCA} (\loc{hierarchical cluster analysis}), es un algoritmo no supervisado que se puede categorizar de dos modos, en aglomerativo o divisivo. El primero funciona con acercamiento \loc{<<bottom-up>>}. Los datos comienzan de forma separada y se van uniendo de forma iterativa según similitud hasta conseguir un único \loc{cluster}. La similitud se puede medir de las siguientes formas:
%         \begin{enumerate}
%           \item Método de \loc{Ward}: La distancia entre dos \loc{clusters} se define por el aumento de la distancia cuadrática después de que los \loc{clusters} sean fusionados.
%           \item Promedio de distancias: Este método usa la distancia media entre dos puntos en cada \loc{cluster}.
%           \item Distancia máxima: La distancia usada es la distancia mayor entre dos puntos en distintos \loc{clusters}.
%           \item Distancia mínima: Se define por la distancia mínima entre dos puntos de dos \loc{clusters} diferentes.
%         \end{enumerate}
%         La medida más común a la hora de calcular estas distancias es la distancia euclídea, aunque otra muy común en la literatura es la distancia \loc{Manhattan}.

%         Por otro lado, el acercamiento divisivo, o \loc{<<top-down>>} funciona dividiendo un único \loc{cluster} basándose en las diferencias entre instancias (puntos del \loc{cluster}). Visualmente se parece a un dendrograma\footnote{Representación gráfica de datos en forma de árbol para organizar datos en categorías y subcategorías hasta alcanzar el nivel de detalle deseado}.

%   \item \textbf{Probabilísticos}: En un algoritmo probabilístico, los puntos se agrupan basándose en la probabilidad de que pertenezcan a una distribución paramétrica \footnote{La estadística paramétrica es una rama de la estadística con la que se deciden valores basándose en distribuciones desconocidas. Estas se determinan según un número finito de parámetros.}. El \sigla{GMM} (\loc{Gaussian Mixture Model}) es el método más usado en \loc{clustereing} probabilístico.
%         \begin{itemize}
%           \item \loc{Gaussian Mixture Model}: Este algoritmo se clasifica como mixto debido a que se compone de un número no especificado de funciones de densidad probabilística. Normalmente se encargará de comprobar si un instancia es parte de una distribución Gaussiana o normal. El algoritmo más usado para determinar la distribución es el \sigla{E-M} (Expectation-Maximizarion) \bib{do2008expectation}.
%         \end{itemize}
% \end{itemize}

% \subsubsection{Reglas de asociación}

% Las reglas de asociación es un método que busca encontrar las relaciones entre variables de un \loc{dataset}. Suelen ser usados en \loc{marketing}, sistemas de recomendación, relación entre productos, estudio de hábitos de consumición \bib{panjaitan2019implementation}, entre otros. El algoritmo más usado en reglas de asociación es el algoritmo <<Apriori>>.

% \begin{itemize}
%   \item \textbf{Algoritmo Apriori}: Este algoritmo permite encontrar <<conjuntos de elementos frecuentes>>. Estos conjuntos sirven para generar reglas de asociación que permiten extender los \loc{sets} a \loc{sets} de mayor tamaño.

%         % \begin{algorithm}[H]
%         %   \caption{Algoritmo Apriori}

%         %   \BlankLine
%         %   \KwIn{Para un \loc{dataset} $T$, un umbral de soporte $\varepsilon$ y un conjunto de datos $C_{k}$ para el nivel $k$}
%         %   \BlankLine

%         %   $L_1 \gets findFrequentOneItemsets(T)$\\
%         %   $k \gets $2\\

%         %   \BlankLine
%         %   \While{$L_{k-1}$ not empty}{
%         %     \tcc{Generación de Apriori}
%         %     $C_k \gets \text{empty list}$\\
%         %     \ForAll{$p \subseteq L\text{, } q \subseteq L $}{
%         %       $c \gets p \cup q_{k-1}$\\

%         %       \ForAll{$u \in L$}{
%         %         \If{$u \subseteq c$}{
%         %           $C_k \gets C_k \cup \lbrace{c}\rbrace$\\
%         %         }
%         %       }
%         %     }

%         %     \BlankLine
%         %     \tcc{Cálculo de candidatos}
%         %     \For{$t \in T$}{
%         %       $D_t \gets \lbrace c \in C_k : c \subseteq t\rbrace$\\
%         %       \For{$c \in D_t$}{
%         %         $count[c] \gets count[c] + 1$\\
%         %       }
%         %     }

%         %     $k \gets k+1$\\
%         %     $L_k \gets \lbrace c \in C_k : count[c] \geq \varepsilon \rbrace $\\
%         %   }

%         %   \BlankLine
%         %   \Return{$Union\left(L_k\right)$}

%         % \end{algorithm}

% \end{itemize}

% \subsubsection{Reducción de la dimensionalidad}

% Normalmente cuando disponemos de una mayor cantidad de datos, tendemos a obtener mejores soluciones, aunque esto tiene impacto en el rendimiento de nuestros algoritmos y puede dificultar comprender de forma completa nuestro \loc{dataset}.
% La reducción de la dimensionalidad se usa cuando el número de características o dimensiones de un \loc{dataset} es muy grande. Reduce el número de instancias a un tamaño manejable manteniendo la integridad de los datos lo máximo posible\bib{sorzano2014survey}\bib{ghodsi2006dimensionality}. Algunos métodos de reducción de la dimensionalidad son: \sigla{PCA} (del inglés \loc{Principal Component Analysis}), \sigla{SVD} (\loc{Singular Value Decomposition}) y \loc{Autoencoders}.

% \subsection{Semi-supervised Learning}

% El aprendizaje semi-supervisado se caracteriza por tener la capacidad de aprender de conjuntos de datos cuyas instancias no están completamente etiquetadas. En nuestro ejemplo \ref{example:car_price}<<cálculo de precio de un coche según sus propiedades>>, tendríamos precios reales solo para algunos de los coches, mientras que otros no tendrían una etiqueta establecida.

% La característica principal que diferencia el aprendizaje semi-supervisado del aprendizaje no supervisado y del aprendizaje supervisado (que veremos a continuación) es que la capacidad de aprender del algoritmo se nutre del hecho de que algunos datos no estén etiquetados. A la hora de intentar averiguar las etiquetas faltantes en base a las presentes, podemos observar comportamientos de aprendizaje que no se verían en los otros dos tipos de \loc{machine learning}.

% De este modo, no es útil solo como algoritmo de aprendizaje si no como algoritmo de generación de nuevos datos en aprendizaje supervisado \bib{zhu2009introduction}. Si la obtención de datos y etiquetas nos supone un proceso costoso y/o el acceso a los datos es limitado \footnote{\pe: La recopilación de datos médicos es algo compleja, ya que se debe tener permiso del paciente, así como hay enfermedades cuyos datos son muy escasos y difíciles de conseguir.}, podemos crear un algoritmo de aprendizaje semi-supervisado para completar la falta de datos.

% Algunos de los métodos usados en aprendizaje semi-supervisado son los siguientes:

% \begin{itemize}
%   \item \textbf{Modelos generativos}: Con este modelo se busca estimar $p(x \mid y)$, la distribución de los puntos pertenecientes a cada clase. La probabilidad $p(x \mid y)$ de que un punto $x$ pertenezca a una etiqueta $y$ es proporcional a $p(x \mid y) \cdot p(y)$ por la Regla de Bayes.

%         Extrayéndolo de \bib{zhu2009introduction}, los modelos generativos asumen que las distribuciones tienen la forma $p(x \mid y, \theta)$ parametrizada por el vector $\theta$. Si esta suposición es incorrecta, los datos no etiquetados pueden incluso disminuir la capacidad de acierto de la solución. Pero, en el caso contrario, si la suposición es correcta, los datos etiquetados mejoraran el rendimiento del modelo.

%         Los datos no etiquetados se suelen distribuir mediante una distribución Gaussiana mixta. La distribución de probabilidad conjunta se puede escribir como $p(x,y \mid \theta) = p(y \mid \theta) \cdot p(x \mid \theta)$ mediante el teorema de probabilidad compuesta. Cada vector $\theta$ se asocia con una función de decisión \[f_{\theta}(x) = \underset{x}{\argmax} p(y \mid x, \theta)\]

%         El parámetro se elige entonces basándonos en la adaptación a los datos etiquetados y no etiquetados, usamos un peso de balanceo $\lambda$:

%         \begin{equation*}
%           \underset{\theta}{\argmax} \left( \log p\left(\left\{x_i, y_i\right\}_{i=1}^{l} \mid \theta\right)+\lambda \log p\left(\left\{x_i\right\}_{i=l+1}^{l+u} \mid \theta \right) \right)
%         \end{equation*}

%   \item \textbf{Modelos basados en grafos}: Este modelo asume un grafo $G = {V,E}$ tal que los vértices $V$ son las instancias de datos (etiquetados y no etiquetados) y $E$ las aristas no dirigidas que conectan las instancias $i,j$ con un peso $w_{ij}$ \bib{song2022graph}\bib{van2020survey}. La construcción de un grafo se cumple en los siguientes pasos. Podemos ver estos pasos en la figura \ref{fig:SSLTaxonomy}:
%         \begin{enumerate}
%           \item \textbf{Construcción del grafo}: Se suele asumir una instanciación inicial del grafo aleatoria
%           \item \textbf{\loc{Weighting}}: Cálculo de pesos de los datos etiquetados para la fase de inferencia.
%           \item \textbf{Inferencia de etiquetas}: La tarea a cumplir aquí es propagar información de las instancias etiquetadas a las <<no etiquetadas>> incorporando la estructura de grafo creada en el paso anterior.
%         \end{enumerate}
% \end{itemize}

% \imagen{./img/memoria/conceptos/Taxonomy_Semi-Supervised_Methods}{Visualización de la taxonomía de los métodos de aprendizaje semi-supervisado. En la rama correspondiente a los métodos basados en grafo podemos observar las distintas fases de clasificación \bib{van2020survey}}{SSLTaxonomy}

\section{Supervised Learning}

También llamado \loc{classification}\footnote{A pesar que no solo trate problemas de clasificación} o \loc{inductive learning} (recordemos fig. \ref{fig:MachineLearningTypes}), el aprendizaje supervisado es aquel aprendizaje cuyas instancias de datos (usadas para que nuestro programa aprenda) están todas etiquetadas.

Según Christopher M. Bishop \et \bib{bishop2006pattern}, el objetivo último del aprendizaje supervisado es ser capaces de observar una variable $t$ (\loc{target}) junto con una variable de entrada $x$ y predecir el valor de $t$ para cualquier nuevo valor de $x$.

Para conseguir la predicción de nuevos datos debemos crear un modelo que va a ser <<entrenado>>, <<testeado>> y <<validado>>. Para conseguir datos para las tres fases dividiremos nuestro \loc{dataset} en tres grupos respectivamente. Con la intención de que el modelo tenga un alto acierto en sus predicciones debe estar bien entrenado, por lo que la mayor fracción de los datos se destinará a esta fase\footnote{Normalmente un 70\% de los datos se asigna a la fase de entrenamiento, aunque depende mucho del problema específico al el que nos enfrentemos}. Pero hay que tener cuidado, porque asignar un porcentaje demasiado alto de los datos a la fase de entrenamiento puede derivar en \loc{overfitting} en el modelo.

El \loc{overfitting} ocurre cuando un modelo se ha sobre-entrenado con unos datos, por que lo que a la entrada de nuevos datos (desconocidos) $x$, generará una predicción $t$ sesgada a los datos de entrenamiento. Para evitar esto, usamos otra fracción de los datos de entrada como fracción <<de control>>: los datos de validación. De este modo iremos comprobando la deriva de las predicciones entre los datos de entrenamiento y los datos de validación.

Ahora podremos controlar la fase de entrenamiento y parar dicho entrenamiento, valga la redundancia, cuando detectemos este sobre-ajuste. Podemos ver esto en la figura \ref{fig:Overfitting}.
Gracias al \loc{set} de validación, detectamos el momento en el que la red comienza a sobre-ajustarse. Esto no solo nos evita cometer errores en la fase de test, sino que mediante la eliminación de iteraciones inútiles sobre el modelo, disminuimos en gran cantidad el tiempo de entrenamiento.

\imagen{./img/memoria/conceptos/Overfitting}{Visualización de \loc{overfitting} en un entrenamiento de una \sigla{CNN} (red neuronal convolucional). Podemos observar en rojo el punto exacto en el que la red comienza a sobre-ajustarse. Este es el momento en el que debemos parar el entrenamiento.}{Overfitting}

En la imagen vemos que el valor del eje de ordenadas tenemos el valor de los costes sobre las iteraciones (\loc{epochs} en inglés). Veremos más sobre esto en el apartado de \loc{fine-tunning} e hiperparametrización de una red neuronal. Veamos ahora los distintos tipos de aprendizaje supervisado.

\subsection{Regression}

En la regresión, el objetivo es predecir un valor numérico y continuo. Uno de los usos principales de modelos de regresión es el relleno de espacios en los datos o la estimación de valores futuros. Por otro lado, la regresión también se puede utilizar para determinar relaciones casuales entre variables dependientes e independientes \bib{maulud2020review}.

En los modelos de regresión, las variables independientes predicen a las variables dependientes. El análisis de regresión estima la variable dependiente $y$ en base al rango de variables independientes $x$. En cuanto a los tipos de regresión, puede ser simple o múltiple.

\begin{itemize}
  \item \textbf{Regresión simple}: En las regresiones simples solo tenemos una variable independiente $y$. Se define la dependencia de la variable como $y = \beta_0 + \beta_1 x + \epsilon$. De forma gráfica, la regresión simple se muestra como una línea entre los puntos cuyo objetivo es minimizar el error cuadrático de las distancias entre dicha línea y los puntos. Es común que en estos modelos tengamos puntos atípicos que se sitúen lejos de la mayoría. Estos puntos se denominan \loc{outliers} y al distanciarse del resto, provocan desviaciones importantes en la regresión \bib{seldon_2021}.
  \item \textbf{Regresión múltiple}: El objetivo de los modelos múltiples es la predicción del conjunto de variables independientes $y$ según el conjunto de variables dependiente $x$. El modelo básico es $y = \beta_0 + \beta_1 x_1 + \cdots + \beta_m x_m + \epsilon$. La fórmula para determinar la matriz es:

        \begin{align*}
          \hat{\beta} & =\left(X^{T} X\right)^{-1} X^{T} y \quad \text { donde }                          \\
          \\
          \quad \beta & =\left[\begin{array}{c}
                                   \beta_{0} \\
                                   \beta_{1} \\
                                   \vdots    \\
                                   \beta_{m}
                                 \end{array}\right], X=\left[\begin{array}{ccccc}
                                                               1      & x_{11}  & x_{12}  & \cdots & x_{1 m} \\
                                                               1      & x_{21}  & x_{22}  & \cdots & x_{2 m} \\
                                                               \vdots & \vdots  & \vdots  & \vdots & \vdots  \\
                                                               1      & x_{n 1} & x_{n 2} & \cdots & x_{m}
                                                             \end{array}\right], Y=\left[\begin{array}{c}
                                                                                           Y_{1}  \\
                                                                                           Y_{2}  \\
                                                                                           \vdots \\
                                                                                           Y_{n}
                                                                                         \end{array}\right]
        \end{align*}

\end{itemize}

\subsection{Classification}

El segundo y último tipo de aprendizaje supervisado es también el más familiar. La clasificación consiste en etiquetar una variable discreta $y$ dada una variable de entrada $x$. Según el número de etiquetas disponibles para la clasificación hablamos de clasificación binaria (2 etiquetas) o multi-clase.

Aunque no listamos todos los métodos, a continuación se pueden observar los más ampliamente usados\bib{ranganathan2018encyclopedia}:

\begin{enumerate}
  \item \textbf{Árboles de decisión}: Los árboles de decisión son una forma de representar reglas sobre los datos de forma jerárquica con una estructura secuencial recursiva que particiona los datos. Los árboles de decisión se usan para aproximar funciones discretas. Un nodo hoja indica el valor de la etiqueta $y$, mientras que un nodo de decisión especifica una operación a realizar. La evaluación comienza evaluando el nodo raíz (un nodo de decisión) y moviéndonos hasta encontrarnos con un nodo hoja, con la etiqueta buscada. El algoritmo más común en la construcción de árboles de decisión es una extensión del algoritmo propuesto por J. Quinlan en 1993 \et \bib{quinlan1993building}: \sigla{ID3} (\loc{Iterative Dichotomiser} 3). El árbol se construye \loc{top-down} de forma recursiva con la técnica de <<divide y vencerás>>. Puedes ver el pseudocódigo a continuación.

        \begin{algorithm}[H]
          \caption{Algoritmo ID3 - TreeBuilding(X, Y)}

          \BlankLine
          \KwIn{Para un \loc{set} de entradas de entrenamiento $X$ y una lista de etiquetas correspondientes $Y$}

          \BlankLine

          $N \gets \text{Nodo raíz}$\\
          \If{$\forall \text{ muestra} = \text{clase } C$}{
            $\text{etiqueta N} = C$\\
            $\Return$\\
          }

          \BlankLine
          \If{$Y = \text{vacío}$}{
            $\text{etiqueta N} = \text{clase más común C en } X \text{(voto mayoritario)}$\\
            $\Return$\\
          }

          \BlankLine
          $y \gets \text{El atributo que mejor clasifica ejemplos donde } y \in Y$\\
          $\text{etiqueta N} \gets y$\\

          \ForAll{$v \in y$}{
            $\text{Haz crecer una rama desde N con condición } y = v$\\
            $X_v \gets \text{\loc{subset} de instancias en Y con } y = v$\\
            \If{$X_v = \text{vacío}$}{
              $\text{añade un nodo hoja etiquetado con la clase más común en X}$\\
            }
            \Else{
              $\text{añade el nodo generado por } TreeBuilding(X_v, Y - v)$\\
            }
          }

        \end{algorithm}

  \item \textbf{\loc{Bagging} y \loc{boosting}}: Es una técnica de ensamblado que combina varios métodos de \loc{machine learning} en un único modelo predictivo en orden de disminuir la varianza (\loc{bagging}) o el \loc{bias(boosting)}.

        El \loc{bagging (Bootstrap Aggregating)} comienza creando muestras aleatorias del \loc{set} de entrenamiento. Tras esto, construye un clasificador para cada muestra. Finalmente, se combinan los resultados de los clasificadores y comienza una votación de mayoría. Al construir estos modelos al mismo tiempo, consideramos a Bagging un método de ensamblado paralelo. Al aumentar el tamaño del \loc{set} de entrenamiento, se disminuye la varianza del modelo final \bib{bauer1999empirical}.

        El \loc{boosting} por otro lado se compone de dos pasos. En el primero usa \loc{subsets} del \loc{set} de entrenamiento para producir modelos con rendimiento medio. Tras esto estimula el rendimiento combinándolos usando el voto mayoritario. En el \loc{boosting}, la creación de \loc{subsets} no es aleatoria y depende del rendimiento de los modelos anteriores.

        La principal diferencia entonces entre los dos, es que \loc{bagging} aprende el \loc{dataset} completo de forma paralela mientars que \loc{boosting} aprende usando los resultados de los modelos anteriores. \loc{Boosting} muestra una mejor \loc{accuracy} que \loc{bagging} \bib{quinlan1996bagging}, pero también tiene a producir \loc{overfitting} en los modelos.

  \item \textbf{\loc{Random Forest}}: El \loc{random forest} también se puede usar en tareas de regresión. El principio que se sigue es que un conjunto de clasificadores <<débiles>> se puede unir para crear un clasificador <<fuerte>>. En este caso tenemos un \loc{ensemble} de árboles de decisión que funciona mejor que \loc{bagging} y \loc{boosting} tanto en \loc{accuracy} como en \loc{overfitting}. Cada árbol del \loc{ensemble} tiene un \loc{set} de entrenamiento distinto. Debemos tener en cuenta que si necesitamos una descripción de las relaciones existentes entre los datos, los \loc{random forest} no son la mejor opción, ya que al tener un conjunto de árboles, se hace difícil la compresión de la estructura general.

  \item \textbf{\loc{Support Vector Machines}} (\sigla{SVM}): Según Van Engelen \et \bib{van2020survey} un \sigla{SVM} es un clasificador que intenta buscar la frontera de decisión que maximice el margen, que a su vez se define como la distancia entre la frontera de decisión y los puntos cercanos a ella\footnote{A veces el margen también se refiere a la zona delimitada por la frontera de decisión en la que caen los puntos}.

        Los \sigla{SVM} también están presentes en la regresión. En el caso de la clasificación, tenemos un conjunto de puntos en un espacio de dimensión $N$. Cada \loc{set} perteneciendo a una de las posibles clases, el \sigla{SVM} busca los planos separados que maximicen los margenes entre los \loc{sets} de datos. A la hora de clasificar un dato, separamos dos planos y calculamos el error cuadrático de los puntos a dichos planos. En algunas ocasiones, los planos pueden no ser linealmente separables; en estos casos, debemos aumentar el número de dimensiones hasta que lo sean. Para esto se usa lo que se denomina función \loc{kernel} \bib{busuttil2003support} \[K(x,y)=\langle f(x), f(y) \rangle \] donde $x$ e $y$ son las entradas de dimensión $n$, y $f$ es un mapa de paso de dimensión-$n$ a dimensión-$m$. Normalmente nos encontraremos con que $m$ es un número mucho mayor a $n$.

  \item \textbf{\loc{Naive Bayes}}: Este clasificador probabilístico se basa en el teorema de Bayes y asume que todas las características están independientemente condicionadas por la etiqueta de la clase \bib{murphy2006naive}. Aunque esto no suele ocurrir, el modelo resultante funciona sorprendentemente bien, a veces compitiendo con técnicas mucho más sofisticadas \bib{rish2001empirical}. Este modelo ha probado ser muy útil en tareas de clasificación de texto y diagnóstico médico, entre otros. En \bib{domingos1997optimality} se prueba que \loc{naive Bayes} funciona de forma óptima en problemas con un alto grado de dependencia entre características.

  \item \textbf{Redes Neuronales}: Por último, tenemos las redes neuronales. Las redes neuronales ocupan la mayoría de la literatura sobre inteligencia artificial y \loc{machine learning} en la actualidad. Una red neuronal es un conjunto de neuronas artificiales colocadas en capas. Cada capa tiene la tarea de extraer características de los datos de entrada, actualizar sus pesos y <<pasar>> información a las siguientes capas.

        Una neurona es una función $f_j$ con una entrada $x = (x_1, \ldots , x_d)$ y un vector de pesos $w_j = (w_{j,1}, \ldots , w_{j,d})$, junto con un \loc{bias} $b_j$ y asociado a una función de activación $\phi$ \bib{bishop1994neural}. La fórmula completa de está función es \[y_j = f_j(x) = \phi (\langle w_j, x \rangle + b_j )\]

        Existen numerosas funciones de activación en una red neuronal, algunas son:
        \begin{itemize}
          \item Función identidad o lineal \[\phi(x) = x\]
          \item \loc{Binary Step} \[\phi(x) = \lbrace\begin{array}{ll}0 & \text{ para } x<0 \\ 1 & \text{ para } x \geq 0\end{array}\]
          \item Función sigmoide \[\phi(x) = \frac{1}{1 + \exp(-x)}\]
          \item Tangente hiperbólica \[\phi(x) = \frac{\exp(2x) -1 }{\exp(2x) + 1}\]
          \item \sigla{ReLU} (\loc{Rectified Linear Unit}) \[\phi(x) = \max(0, x)\]
          \item \loc{Leaky} \sigla{ReLU} \[\phi(x) = max(0.1 \cdot x,x)\]
          \item \sigla{ELU} (\loc{Exponential Linear Unit}) \[\phi(x) = \lbrace\begin{array}{ll}x & \text{ para } x \geq 0 \\ \alpha(\exp(x) - 1) & \text{ para } x < 0\end{array}\]
          \item Función softmax: descrita como una combinación de múltiples sigmoides, calculamos su salida con \[\phi(x_i) = \frac{exp(x_i)}{\sum_j exp(x_j)}\] siendo $x$ el conjunto de las salidas de una neurona en un problema multiclase.
        \end{itemize}

        Tanto la función de \loc{binary step} como la función de identidad (ambas funciones lineales) no permiten \loc{backpropagation} ya que la derivada de la función es una constante y no tiene relación alguna con la variable de entrada x (no aporta datos adicionales) \bib{hinton2006unsupervised}.

        Por otro lado, en regresión se usan funciones de activación lineales, mientras que usaríamos sigmoides en clasificación binaria, o softmax en clasificación multietiqueta.

        Una vez tenemos nuestra función de activación elegida, podemos calcular la salida de la red dada una variable de entrada $x$. Y una vez calculada la salida de la red, podemos comparar el resultado $y$ con la etiqueta de \loc{ground truth} correspondiente con dicha entrada. Según sea igual o no, modificaremos el vector de pesos de nuestra neurona.

        Repetiremos este proceso hasta que el \loc{accuracy} de la neurona (o la red de neuronas) sea superior al \loc{threshold} buscado \bib{ranganathan2018encyclopedia}.

        Podemos medir el \loc{accuracy} de una neurona dividiendo el número de predicciones correctas entre el número total de predicciones \bib{flach2019performance}. Pero el \loc{accuracy} no es el único método de medida de acierto en las redes neuronales. También tenemos la precisión, el \loc{recall}, el \loc{F-measure}, la especificidad, el \loc{fallout} o la Coeficiente de Correlación de \loc{Mathew} (\sigla{MCC}) \bib{chicco2020advantages}.

        Con
        \begin{align*}
          tp = \loc{true positives}  \\
          tn = \loc{true negatives}  \\
          fp = \loc{false positives} \\
          fn = \loc{false negatives}
        \end{align*}

        \begin{itemize}
          \item \loc{Accuracy}: $Acc = \frac{tp + tn}{tp + tn + fp + fn}$
          \item Precisión: $P = \frac{tp}{tp + fp}$
          \item \loc{Recall}: $R = \frac{tp}{tp+ fn}$
          \item \loc{F-measure}: $F = \frac{2P \cdot R}{P + R}$
          \item Especificidad: $S_p = \frac{tn}{tn + fp}$
          \item \loc{Fallout}: $\frac{fp}{tn + fp}$
          \item Coeficiente de Correlación de \loc{Mathew}: $\sigla{MCC} = \frac{tp \cdot tn - fp \cdot fn}{\sqrt{(tp + fp)(tn + fn)(tp + fn)(tn + fp)}}$
        \end{itemize}
\end{enumerate}

\section{Deep Learning}

Hemos visto en el último apartado las características de una red neuronal de forma individual: como se compone, que tipos de funciones de activación tenemos y como podemos evaluar y medir la salida de una neurona.

Las redes neuronales más básicas: \sigla{FFNN} (\loc{Feed Forward Neural Netowrks}) y Perceptrones, se componen únicamente de una capa de entrada y salida, con ocasionalmente una o más capas ocultas entre ambas. El método para entrenar estas redes suele ser el \loc{backpropagation}.

El término \loc{backpropagation} viene tras ser empleado por \loc{Rosenblatt} (1962) \bib{block1962analysis} a la hora de intentar generalizar el entrenamiento de un perceptrón a múltiples capas ocultas. No fué muy exitoso hasta que \loc{Rumelhart} y \loc{Williams} (1986) \bib{rumelhart_hinton_williams_1986} extendieron la idea y la hicieron popular entre investigadores \bib{rumelhart1995backpropagation}. La idea del \loc{backpropagation} es minimizar el error entre la salida correcta e incorrecta mediante el \loc{gradient descent}. Para esto necesitamos calcular las derivadas parciales de nuestras funciones de activación \footnote{Recordemos que solo las funciones de activación no lineales pueden ser usadas para el \loc{backpropagation}, ya que la derivada de una función lineal es una constante}. Calculamos la salida de la derivada en el \loc{forward pass}, y con el resultado obtenido, realizamos un \loc{backward pass}, con el que propagamos y actualizamos los pesos de la red.

Esto ahora nos parece básico, pero fue una <<revolución>> en su momento y es la base principal del \loc{deep learning}. Pero, el \loc{deep learning} no solo trata el aprendizaje supervisado, (también tenemos \loc{deep reinforcement learning}, \loc{unsupervised learning}\footnote{Una de las redes neuronales en \loc{unsupervised learning} más habladas en la literatura son las redes \sigla{GAN} (\loc{Generative Adversarial Neural Networks}). Estas redes usan la estructura de discriminador-generador \bib{GenerativeAdversarial} para extraer información de los datos de entrada} y \loc{semi-supervised learning}). Por simplificar vamos a tratar solo redes neuronales de aprendizaje supervisado o \sigla{DNN} (\loc{Deep Neural Networks}). Una lista extensa de redes neuronales se puede ver en \bib{typesann}.

Las \loc{deep neural networks} (\sigla{DNN}) se clasifican en dos tipos principales: Redes Neuronales Recurrentes (\sigla{RNN}) y Redes Neuronales Convolucionales (\sigla{CNN}). Se diferencia de las redes neuronales artificiales <<clásicas>> (\sigla{ANR}) \footnote{Como \sigla{FNNN} o perceptrón} en que en vez de tener una capa de entrada seguida de una o más capas escondidas, con una capa de salida a continuación; las \sigla{DNN} tienen más de una capa entre la salida y la entrada, de ahí que se denominen profundas \bib{typesann}.

\imagen{./img/memoria/conceptos/RNN_CNN}{Estructura de una red \sigla{RNN} y \sigla{CNN}}{RNN_CNN}

En la figura \ref{fig:RNN_CNN}, se puede observar la estructura de los dos tipos mencionados. La red recurrente debe su nombre a las neuronas que están conectadas consigo mismas (neuronas en azul).

Las redes neuronales convolucionales son más complejas y su particularidad, como veremos más en profundidad en el siguiente apartado, es la extracción de características y reducción de la dimensionalidad. Vemos que la capa de entrada (en amarillo) tiene 5 neuronas, y que en las siguientes capas, el número de neuronas se va reduciendo hasta llegar a dos en la capa justo anterior a la oculta (en verde). Esto se produce mediante la convolución o el \loc{pooling}. Vamos a estudiar estas dos redes y sus términos con un poco más de profundidad a continuación.

\begin{itemize}
  \item \textbf{\sigla{RNN}}: Las Redes Recurrentes son redes que tuvieron su auge en investigación en los años 90. \loc{Fausett} (1994) \bib{fausett1994fundamentals} las define como redes neuronales con conexiones \loc{feedback} (retroalimentación). Ejemplos de estas redes son las redes \sigla{BAM} (\loc{Bidirectional associative memory}), redes \loc{Hopfield} y máquinas de \loc{Boltzmann} \bib{medsker1999recurrent}.

        La característica principal de las redes recurrentes es que tienen un estado que mantienen entre distintas iteraciones de la red. Las neuronas no solo reciben información de la capa anterior si no que también reciben información de si mismas y de la iteración anterior \bib{elman1990finding}. Por esta razón, el orden con el que <<alimentamos>> datos a la red es importante y puede cambiar la salida generada de forma radical. Esta característica las hace muy útiles con formatos de datos como vídeo o audio \footnote{Aunque también podemos tratar texto o imagen como concatenación de letras o píxeles en el tiempo}. Su principal aplicación es el avance y completado de información, como por ejemplo el autocompletado de texto.  En \bib{rodriguez1999recurrent} podemos ver un ejemplo de una \sigla{RNN} aplicada para enseñar a la red a contar.

  \item \textbf{\sigla{CNN}}: Y llegamos a las redes convolucionales. Introducidas por \loc{Le Cun} (1998) \bib{le1989handwritten}, estas redes son uno de los tópicos más hablados en la literatura actual en el mundo del \loc{machine learning} y la inteligencia artificial. Estas redes se usan en muchísimas aplicaciones y dominios, sobre todo en procesamiento de vídeo e imagen.

        Las redes convolucionales tienen una estructura muy diferente a cualquier otra red neuronal (figura \ref{fig:CNN_Arquitecture}). Se componen de capas de convolución, capas de \loc{subsampling} y capas <<\loc{fully-connected}>>.

        \imagen{./img/memoria/conceptos/CNN_Arquitecture}{Estructura de una \sigla{CNN}. Figura de \loc{Le Cun} \et \bib{lecun1998gradient}}{CNN_Arquitecture}

        A continuación vamos a hablar de filtros o \loc{kernels}, \loc{pooling}, \loc{padding}, \loc{stride}, \loc{ResNet}, entre otros términos.

        \begin{enumerate}
          \item \textbf{Extracción de características o convolución}: En esta fase se busca reducir la dimensionalidad de la entrada manteniendo las características espaciales de los datos. De esta forma no perdemos información y es más fácil procesarla en las siguiente capas. En redes neuronales tradicionales, transformaríamos una entrada (imagen o vídeo) en un vector de dimensión uno.

                Al hacer esto perderíamos información de formas en la imagen, o de movimientos de objetos o personas en los vídeos. Por esa razón, la extracción de características de una imagen se hacía antes de forma manual en la fase de preparación de los datos.

                Para eliminar esta fase manual, las \sigla{CNN} actúan directamente en matrices de datos multidimensionales (o tensores \footnote{Un tensor es un objeto algebraico que representa un espacio vectorial. Puede mapear vectores y matrices multidimensionales \bib{wiki:tensor}}), en vez de actuar sobre vectores. Pero, ¿cómo extraemos información de un tensor?

                Para extraer información de un tensor manteniendo sus características espaciales vamos a usar \loc{kernels} o filtros.

                Un filtro es un tensor con las mismas dimensiones que el tensor de entrada de la red, pero con tamaños mucho menores. \pe: si tenemos como entrada una imagen de $(200x200)$ píxeles, tendremos un filtro de dos dimensiones, por ejemplo, de $(3x3)$ píxeles. Si en vez de una imagen, estamos clasificando un vídeo con 10 \loc{frames} $(10x200x200)$ tendremos un filtro que por ejemplo puede tener el tamaño $(2x3x3)$, siendo 2 el número de \loc{frames}.

                Una vez tenemos nuestro tensor de entrada y nuestro \loc{kernel}, podemos iniciar el proceso de convolución. Para hacer esto, vamos a <<mover>> nuestro \loc{kernel} por nuestro tensor de entrada, generando un <<mapa de características>>. Veamos un ejemplo:

                \imagen{./img/memoria/conceptos/Kernel}{Ejemplo de convolución con \loc{kernel} $3x3$ sin \loc{strading} ni \loc{padding}}{Kernel}

                Vemos en la figura \ref{fig:Kernel} que tenemos un tensor bidimensional (\pe una imagen) de tamaño $(5x5)$ con un \loc{kernel} de tamaño $(3x3)$. El mapa de características resultante tiene un tamaño de $(3x3)$\footnote{Da la casualidad de que el tamaño de la salida es igual al tamaño del kernel, pero no es la norma. Si tuviéramos una entrada de $(5x7)$ obtendríamos una salida de $(3x5)$}.

                En este caso no hemos tenido ningún problema y hemos podido <<encajar>> nuestro \loc{kernel} en la matriz perfectamente. Esto no siempre es así, ya que si queremos adaptarnos a cualquier tipo de entrada, es posible que tengamos tamaños cuyo \loc{kernel} no <<encaje>> de forma perfecta.

                Por ejemplo, si nuestra matriz fuera de $(4x4)$ en vez de $(5x5)$, tendríamos que para el tercer cálculo, nuestro \loc{kernel} se sitúa fuera de la matriz (en el lado derecho). Para solucionar este problema usamos \loc{padding} o \loc{stride}.

                \begin{itemize}
                  \item El \loc{padding} se usa para añadir <<información>> en los bordes de nuestro tensor y así mantener el tamaño de salida deseado. Podemos elegir añadirlo en todos los bordes o solo en algunos. El valor óptimo del \loc{padding} es $0$. Esto es porque toda información que añadamos a los bordes, <<contamina>> la información importante ubicada en el tensor de entrada \bib{o2015introduction}.

                  \item El \loc{stride} por otro lado indica la cantidad de unidades que movemos nuestro \loc{kernel} sobre nuestra entrada en cada iteración de la convolución. Valores comunes de \loc{stride} son $2$ o $3$ \bib{gu2018recent}. Tenemos que tener cuidado aquí, porque con un \loc{stride} demasiado alto, podemos saltarnos extracciones de características importantes de nuestro tensor.
                \end{itemize}

                Podemos calcular el tamaño de salida con la siguiente fórmula \bib{o2015introduction}: \[\delta = \frac{D + 2p - K}{s} + 1\] donde $D$ es un vector de tamaño $d$ con las dimensiones de nuestro tensor de entrada, $K$ es un vector de tamaño $d$ con las dimensiones del \loc{kernel}, y $p$ y $s$ son los valores de \loc{padding} y \loc{stride} respectivamente. La salida será un número $\delta$. Si el número es un entero, entonces la convolución será correcta.

                Así en nuestro caso $\delta = [\frac{5 + 2 \cdot 0 - 3}{1} + 1, \frac{5 + 2 \cdot 0 - 3}{1} + 1] = [3, 3]$. Como vemos, en el proceso de convolución, <<alteramos>> las dimensiones del tensor de entrada con la finalidad de extraer las características espaciales del mismo.

          \item \textbf{\loc{Pooling}}: Muy similar a la fase de convolución, el \loc{pooling} reduce el tamaño del mapa de características. El objetivo es disminuir la capacidad computacional necesaria para procesar gran cantidades de datos \bib{cui2017kernel}. Hay dos tipos de \loc{pooling} principales: \loc{Max Pooling} y \loc{Average Pooling}. El primero devuelve el valor máximo de la <<porción>> del tensor cubierta por el \loc{kernel} mientras que el segundo devuelve la media de todos los valores en dicha <<porción>>. \loc{Max pooling} funciona normalmente mejor que \loc{average pooling} porque no solo lleva a cabo la labor de reducción de la dimensionalidad, sino que también actúa como mecanismo de eliminación de ruido. \footnote{Denotar que lo importante del \loc{kernel} en el \loc{pooling} son sus dimensiones ya que no vamos a realizar operaciones con él}.

          \item \textbf{Capa \loc{fully connected}}: Una vez hemos extraído las características necesarias de nuestro tensor de datos, hemos terminado con un tensor con dimensiones reducidas. Esta fase se ve muy bien gráficamente en la figura \ref{fig:RNN_CNN}. Las capas de color rosa son las capas en las que realizamos la convolución y el \loc{pooling}. El tensor resultante debe ahora sí, ser convertido en un vector (tensor de dimensión uno), ya que este es la entrada ahora de una capa \loc{fully connected}. Este proceso se denomina \loc{flatten}.

                Las capas \loc{fully connected} se asemejan mucho a estructuras tradicionales de una red neuronal, con la diferencia de que las neuronas que forman esta capa aplican una transformación lineal del vector de entrada primero, seguido de una transformación <<no lineal>>.

                La salida de esta capa se calcula como \[
                  y_{j k}(x)=f\left(\sum_{i=1}^{n} w_{j k} x_{i}+w_{j 0}\right)
                \] donde $w_0$ es el \loc{bias}, $w$ nuestra matriz de pesos, $x$ el vector de entrada y $f$ la función de activación no lineal.

                El nombre de esta capa \footnote{También se suele llamar capa densa} se debe a que las neuronas están totalmente conectadas entre si entre las distintas capas (no entre neuronas de la misma capa ni consigo mismas) \bib{huang2017densely}.

                La salida de la capa \loc{fully conneted} será la salida que produce la función de activación no lineal que se aplica tras las funciones lineales. Esta función puede ser una función \loc{softmax} para clasificación multiclase o puede ser una sigmoide para otro tipo de salida. Por ejemplo, en este proyecto, se ha usado una \loc{fully connected} con \loc{softmax} para clasificar los signos y una \loc{fully connected} con sigmoide para extraer un vector de poses de gestos en un vídeo. Puedes avanzar a la sección <<Aspectos relevantes>> para más información.
        \end{enumerate}

        Existen numerosas arquitecturas de redes neuronales convolucionales que han demostrado ser muy útiles en la resolución de distintos problemas. Algunas son: \sigla{AlexNet} \bib{krizhevsky_sutskever_hinton_2012}, \sigla{LeNet}\bib{lenet}, \sigla{VGGNet} \bib{VGGNet}, \sigla{GoogLeNet} \bib{GoogLeNet}, \sigla{ZFNet} \bib{ZFNet} o \sigla{ResNet} \bib{ResNet}.
\end{itemize}

Todas estas no son ni mucho menos todas las redes neuronales del campo del \loc{deep learning}. Cada pocos meses se publica un \loc{paper} que muestra estructuras nuevas con oportunidades futuras impresionantes. En la actualidad, otro de los modelos de \loc{deep learning} que más atención (nunca mejor dicho) genera, son los modelos \loc{transformers}, expuestos por primera vez en el famoso \loc{paper} <<\loc{Attention is all you need}>>, parte de un libro de (2017) \bib{attention}. Estos modelos usan el llamado <<mecanismo de atención>> que da pesos diferentes a las distintas partes de nuestro dato de entrada (llamados \loc{tokens}). Los principales campos en los que se usan los \loc{transformers} son el campo del \loc{natural language processing} (\sigla{NLP}) y el campo de \loc{computer vision} (\sigla{CV}).

\textbf{Nota del autor}: Los \loc{transformrs} se alejan un poco del tema principal de este proyecto, pero invito a leer el fantástico \loc{paper} y descubrir sobre los modelos que usan esta red y que están revolucionando el mundo del \loc{deep learning} (\pe: \sigla{GPT-3} o \sigla{DALL-E})

\section{Data Augmentation}

Las redes neuronales convolucionales, como hemos visto en los apartados anteriores, tienden a funcionar
extraordinariamente bien en tareas de clasificación de vídeo e imagen, pero, en contrapartida, necesitan gran cantidad de datos para evitar \loc{overfitting} \bib{shorten2019survey}.

Desgraciadamente, son muchas las ocasiones en las que no podemos usar datos de entrenamiento suficientes para evitar \loc{overfitting}, como en el caso de datos médicos. En estos casos, aplicamos aumento de datos.

\loc{Data augmentation} o aumento de los datos, agrupa un conjunto de técnicas con las cuales buscamos extender el conjunto de datos de entrenamiento de nuestros modelos, en caso de que sean insuficientes.

Las principales técnicas de aumento de datos son la transformación geométrica, las transformaciones en el espacio de color, aplicación de filtros \loc{kernel} (como hemos visto en la convolución), mezcla de distintas imágenes o vídeos, borrado aleatorio, aumento del espacio de características, transferencia de estilo y redes generativas adversarias (\sigla{GANs}).

El aumento de datos busca un acercamiento al \loc{overfitting} desde la raíz principal del problema: el \loc{dataset}. Esto es asumiendo que el \loc{dataset} tienen datos que pueden ser aumentados.

Estas aumentaciones artificiales inflan el \loc{dataset} en tamaño de dos formas: por deformación de datos (\loc{data warping}) o por \loc{oversampling} (sobremuestreo).

Las deformaciones de datos transforman los datos existentes preservando su etiqueta. Algunos tipos son:

\begin{enumerate}
  \item Transformación geométrica y/o de color
  \item Eliminación parcial aleatoria de datos espaciales
  \item Transferencia de estilo (mediante redes adversarias generativas (\sigla{GANs})). Podemos ver un ejemplo en la figura \ref{fig:CycleGANs}.
\end{enumerate}

Por otro lado, el supersampleo consiste en mantener los datos actuales y añadir nuevas instancias al \loc{set} de entrenamiento. La aplicación de \loc{oversampling} y \loc{data warping} no es una dicotomía mutuamente exclusiva.

\imagen{./img/memoria/conceptos/dataAugmentation}{Ejemplo de \loc{data augmentation} aplicando \loc{adversarial neural netowkrs}, más concretamente \loc{CycleGANs}. Extraído de \bib{face_CycleGANs}}{CycleGANs}

Lo que hace el reconocimiento de imágenes y vídeo un problema tan complejo en \loc{deep learning} es que el modelo de reconocimiento debe sobrellevar problemas del punto de visión, la luz, la oclusión ambiental, el color del fondo, la escala de los elementos, entre otros. El aumento de datos consiste en internalizar estas pequeñas variaciones de movimiento y color en los datos, de tal modo que, una vez entrenado el modelo, funcione en todas las situaciones posibles.

Pero el \loc{data augmentation} no es la única técnica que se puede aplicar para evitar \loc{overfitting}. Otra técnica bastante común (y que nosotros hemos aplicado) es el denominado \loc{batch normalization}. El \loc{batch normalization} es una técnica de regularización de los datos que normaliza los \loc{sets} de las capas de activación de una red neuronal. Esta normalización funciona restando la media de cada activación y dividiendo entre la desviación estándar de todos los valores del \loc{batch} \bib{shorten2019survey}.

\capitulo{4}{Técnicas y herramientas}

% Esta parte de la memoria tiene como objetivo presentar las técnicas metodológicas y las herramientas de desarrollo que se han utilizado para llevar a cabo el proyecto. Si se han estudiado diferentes alternativas de metodologías, herramientas, bibliotecas se puede hacer un resumen de los aspectos más destacados de cada alternativa, incluyendo comparativas entre las distintas opciones y una justificación de las elecciones realizadas. 
% No se pretende que este apartado se convierta en un capítulo de un libro dedicado a cada una de las alternativas, sino comentar los aspectos más destacados de cada opción, con un repaso somero a los fundamentos esenciales y referencias bibliográficas para que el lector pueda ampliar su conocimiento sobre el tema.



\capitulo{5}{Aspectos relevantes del desarrollo del proyecto}

En esta sección vamos a hablar de algunos de los problemas y retos con los que nos hemos enfrentado en el proyecto. Vamos a hablar de la carga de los datos y de como pasamos de clasificaciónd de imágenes a clasificación de video. Seguiremos hablando un poco del tratamiento de los datos. Tras esto mencionaremos una de las estructuras de \sigla{CNN} usada, ResNet. Continuamos hablando de la estructura de nuestra red y problemas que hemos tenido, junto con la hiperparametrización y el entrenamiento. Para terminar hablaremos de como hemos usado una red con dos salidas para intentar mejorar el \loc{accuracy} y de como se ha explorado la cuantización para comprimir el modelo y hacerlo <<\loc{deployment-ready}>>.

\section{Carga de datos}

El objetivo de este proyecto es la clasificación de lenguajes de signo a nivel palabra. El \loc{dataset} usado \bib{li2020word} es \sigla{WLASL} (Word-Level American Sign Language) y está extraido de un paper que ganó mención honorable en el \sigla{WACV} (2020). Este \loc{dataset} se compone de 21083 videos, que se clasifican en 2000 palabras distintas.

Al comienzo del proyecto, pensabamos que usando \prog{PyTorch} ibamos a tener las herramientas necesarias para cargar los datos en formato video y procesar las transformaciones directamente. Esto resultó no ser así.

Nos dimos cuenta que no teniamos una forma directa de cargar un video con \prog{Pytorch} y transformarlo a un tensor con el en el cual podíamos ejecutar transformaciones o usar como entrada.

Por eso, lo primero que hicimos, fue crear unos \loc{scripts} con los que transformar los videos en \loc{frames}. Con la idea de poder ser reutilizado lo máximo posible, nos pusimos a diseñar los \loc{scripts} de forma dinámica, es decir; los \loc{scripts} debían aceptar distintos argumentos para procesar distintos \loc{datasets}, independientemente de que el más importante fuera el citado arriba.

Puedes ver más con detalle la estructura del script en el anexo <<Manual del programador>> pero básicamente el \loc{script} se puede ejecutar con las siguiente opciones:

GLOBALES
\begin{itemize}
  \item -i, --input: Carpeta con los videos a procesar (Opcional).
  \item -o, --output: Carpeta donde se guardaran los frames (Opcional).
  \item -l, --labels: Número de labels a procesar. Puede ser un número o ``all`` para procesar todos (Opcional).
  \item -c, --config: Archivo de configuración con las etiquetas (Opcional).
  \item -h, --help: Muestra esta ayuda.
\end{itemize}

EXTRACCIÓN DE FRAMES
\begin{itemize}
  \item -m, --merge: Indica si se debe unir los frames en una imagen (Opcional).
  \item -f, --frames: Número de frames aleatorios del video (Opcional).
\end{itemize}

EXTRACCIÓN DE VIDEOS
\begin{itemize}
  \item -v, --video: Extrae todos los frames del video en una carpeta única para ese video. Deshabilita las opciones -f y -m (Opcional).
\end{itemize}

Con estas opciones nos ahorramos tener que estar transformando los datos de forma <<manual>> o haciendo cambios a los \loc{scripts} cada vez que queríamos usar unas u otras etiquetas.

Bien, una vez nos creamos estos \loc{scripts} para transformar de video a \loc{frames}, estabamos listos para cargarlos con \prog{PyTorch} en un \prog{DataLoader} (clase de \prog{PyTorch} usada para cargar datos desde la estructura de archivos).

Como funciona el \prog{DataLoader}, usa la estructura de archivos para cargar los datos, es decir, que carga los datos y los relaciona con la etiqueta correspondiente según el nombre de la carpeta en la que están.

\imagen{./img/memoria/aspectos/FolderStructure}{Estructura de carpetas necesaria para la carga de datos con \prog{PyTorch}}{FolderStructure}

En la figura \ref{fig:FolderStructure} podemos observar que los \loc{frames} exportados se situan en carpetas según la etiqueta que los clasifica. Al tenerlo de este modo, podemos cargar todo el dataset con \prog{PyTorch} y hacer dentro la división en \loc{set} de entrenamiento, test y validación.

Nota: Esta estructura de clases se mantiene para todos los dataset. La figura que acabamos de ver se corresponde con la red que clasifica únicamente \loc{frames} y no videos (la primera red construida). La estructura de carpetas para el dataset con videos es muy similar, solo que en vez de haber un frame por video en cada clase, tenemos una carpeta por video con los frames de ese video, para cada clase (un nivel más).

\subsection{Volvemos a unir los frames}

Una vez tenemos la estructura lista para cargar los datos con \prog{PyTorch}, debemos unir los frames en un tensor <<video>> (en el caso de que tengamos más de un frame por video). Para esto usamos un \loc{DataLoader} personalizado, que va a concatenar las imagenes en un tensor único. De este modo pasamos de tener un tensor de tamaño $(3xHxW)$ a uno de tamaño $(Fx3xHxW)$ donde la $H$ y la $W$ representan la altura y la anchura de un \loc{frame} respectivamente, y la $F$ el número de frames cargados para un video en particular \footnote{El 3 identifica los tres canales \sigla{RGB} de una imagen en color}.

Una vez hemos hecho la concatenación, ya tenemos nuestro dataset tarnsformado en tensores de 4 dimensiones. Ahora debemos dividir los datos en \loc{sets} de entrenamiento, test y validación, no sin antes, hacer las transformaciones necesarias en los datos y agrupar los videos en \loc{batches}.

%%%%%%%%%%%%%%%%%%%%%%%%%%%%%%%%%%%%%%%%%%%%%%%%%%%%%%%%%%%%%

\section{Tratamiento de los datos}

Una vez hemos cargado los datos, comienza una de las parte más importantes en el \loc{machine learning} y la minería de datos: la transformación de los datos y adaptación para el entrenamiento.

En nuestro caso, al usar \sigla{CCNs}, y como hemos visto en el apartado teórico, el tratamiento de datos no es tán importante. Esto se debe a que en la fase de convolución, el \loc{kernel} tendrá en cuenta todas las dimensiones presentes en nuestro tensor de entrada, y se entrenará para sacar las caracterísitcas que más importen.

Aún así, hay una transformación en los datos que es muy importante, ya que debido al alto consumo de las \sigla{CNNs}, necesitamos un ajuste sobre el tamaño de las imagenes, pasando del tamaño original a $(224x224)$ (píxeles).

Una vez hecha esta transformación sobre las imágenes, normalizamos los datos.

\subsection{Normalización}

La normalización, que consiste en trasladar valores de su rango de valores original a otro controlado (normalmente entre 0 y 1), es muy útil cuando estamos tratando entardas con datos que se situan en rangos muy distintos \footnote{\pe: un valor enumerado entre 0 y 5 y un valor continuo entre 100 y 200}.

En nuestro caso la entrada de los datos está compuesta únicamente por imágenes de color (esto son 3 canales con pixeles que varian entre 0 y 255), por lo que al no mezclar con otro tipos de datos y rangos, podríamos proceder sin normalizar.

Pero no es la mejor opción. normalizar los datos es una práctica común para conseguir acercar la media de los datos a 0 y acelera el aprendizaje de la red proporcionando una convergencia más rápida.

\subsection{\loc{Data Augmentation}}

El \loc{data augmentation} consiste en aumentar el tamaño del \loc{dataset} duplicando datos del mismo y alterándolos para que parezcan distintos.

En nuestro caso no hemos incluido esta técnica debido a tratar con videos y no con imágenes. Al tener la transformación de los datos a nivel imágen (\loc{resize} y normalización), se nos complicaba la aumentación de los datos ya que debiamos <<alterar>> los datos para todos los \loc{frames} del mismo video, con tal de no alterar la información espacial de dicho video.

Esto es algo que nos gustaría marcar para un plan futuro en el proyecto, ya que el \loc{dataset}, aunque muy extenso, tiene muy pocos videos por etiqueta (aproximadamente 40), y aumentar este número aunque sea en 10, sería una extensión notable del dataset.

\subsection{\loc{Batching}}

A la hora de entrenar una red, no podemos usar el \loc{dataset} completo como entrada, ya que nos quedaríamos sin memoria enseguida. Por esto agrupamos las entradas en \loc{batches} o partes. Al hacer esto, tenemos que modificar nuestra red para aceptar los datos en grupos.

El \loc{batching} realmente lo que está haciendo es aumentar la dimensión de nuestro tensor de entrada. En los puntos anteriores hemos visto que nuestro tensor de entrada era del tipo $(Fx3xHxW)$ y que representaba un video de nuestro dataset. Ahora añadimos una dimensión más $B$ representando el tamaño de cada \loc{batch}: $(BxFx3xHxW)$, resultando en un tensor de entrada de 5 dimensiones.

Cuando entrenamos una \sigla{CNN} con un tensor de entrada de 5 dimensiones decimos que tenemos una red convolución 3D \footnote{Una \sigla{CNN} 1D se usa para procesar datos lineales, como la altura según el tiempo. Una \sigla{CNN} 2D se usa para imágenes y una 3D (que significa de 3 o más dimensiones) para videos u otros datos que necesitan imágenes en \loc{stack} }.

\subsection{\loc{Subsampling}}

En el subsampling divisimos los \loc{batches} creados en \loc{sets} de entrenamiento, test y validación.

En nuestro caso hemos usado un \loc{subsampler} aleatorio porporcionado por \prog{PyTorch}: \prog{SubsetRandomSampler} con una división de 70\%, 20\% y 10\% para entranamiento, test y validación.

%%%%%%%%%%%%%%%%%%%%%%%%%%%%%%%%%%%%%%%%%%%%%%%%%%%%%%%%%%%%%

\section{Red Neuronal ResNet}

Antes de comenzar a construir nuestra \sigla{CNN}, probamos el \loc{dataset} con una \loc{ResNet} ya construida. Una ResNet \bib{ResNet} es una estructura que usa numerosas capas de \sigla{CNNs} una detrás de la otra y conectadas entre sí cada dos o tres capas (conexión residual). Tenemos varios tipos identificadas por el número de capas que tienen: \loc{ResNet18}, \loc{ResNet34} \loc{ResNet50}, \loc{ResNet101}, etc.

En nuestro caso, la prueba que hicimos con una \loc{ResNet} fue usando un dataset con fotos de animales \bib{alessio_2019} con 10 etiquetas y con aproximadamente 50 ejemplos por etiqueta. El accuracy conseguido en entrenamiento y test (ya que en este punto incial no teníamos validación) es:

\tablaSmall{Accuracy de una red ResNet con un dataset pequeño de imágenes de animales \bib{alessio_2019}}
{l c c}{acc_resnet18}
{ Tipo de ResNet & Entrenamiento & Test \\}{
  ResNet18 & 97,59\% & 95,51\%  \\
}

Por otro lado, la misma red con nuestro dataset compuesto por un único frame por video consiguió el siguiente \loc{accuracy}:

\tablaSmall{Accuracy de una red ResNet con nuestro dataset original con 10 etiquetas, 30 videos por etiqueta, 1 solo frame por video}
{l c c}{acc_resnet_frame}
{ Tipo de ResNet & Entrenamiento & Test \\}{
  ResNet18 & 92,4\% & 64,6\%  \\
}

Como vemos, mucho menor. Pero esto es normal, ya que clasificar un signo completo de un video con un solo frame es imposible. A la hora de probar el modelo con una inferencia por webcam, los resultados eran prácticamente aleatorios.

Al ver esto, decidimos entonces pasar ya a crear nuestra propia red convolucional.

%%%%%%%%%%%%%%%%%%%%%%%%%%%%%%%%%%%%%%%%%%%%%%%%%%%%%%%%%%%%%

\section{Estructuras de la redes convolucionales}

La estructura de nuestra primera red convolucional (denominada <<SimpleNet>>), usada para clasificación de video, es mucho más simple que la que vemos en una red \loc{ResNet}. Se compone de tres capas de convolución con \loc{pooling} y una capa densa (o \loc{fully connected}) con activación \loc{softmax} al final. Podemos ver la estructura en la figura \ref{fig:1stCNN}.

\imagen{./img/memoria/aspectos/1stCNN}{Estructura de una \sigla{CNN} exportada con \loc{Netron}}{1stCNN}

La segunda red convolucional que tenemos, que hemos nombrado <<TwoOutNet>>, tiene una estructura que comienza con dos capas convolucionales. Tras esto, tenemos una gran diferencia con la primera, y es que en este caso, generamos dos salidas (la clase y la pose detectada). Por eso tenemos dos capas densas, una por salida, con una función de activación de la capa oculta \loc{softmax} y \loc{sigmoid} respectivamente.

Podemos ver en la figura \ref{fig:2ndCNN} la estructura mencionada con la división en dos tras la segunda convolucional.

\imagen{./img/memoria/aspectos/2ndCNN}{Estructura de una \sigla{CNN} con doble salida exportada con \loc{Netron}}{2ndCNN}

Para saber más sobre el porque de la doble capa de salida, adelanta al apartado de <<Dos salidas en la misma red>>

%%%%%%%%%%%%%%%%%%%%%%%%%%%%%%%%%%%%%%%%%%%%%%%%%%%%%%%%%%%%%

\section{Entrenamiento e hiperparametrización}

Los hiperparámetros que hemos cambiado en las pruebas son:

\begin{itemize}
  \item Número de segmentos: Esto indica el número de partes en los que dividimos un video.

  \item \loc{Frames} por segmento: Con este valor cambiamos el número de frames que seleccionamos en cada segmento.

  \item \loc{Batch size}: El tamaño del \loc{batch} de entrada de la red. Un número muy alto necesita mucha memoria de \sigla{GPU}.

  \item Tamaño de la imagen de entrada: En píxeles, aquí indicamos el tamaño que se debe usar en la fase de transformación de los datos antes de entrar en la red.

  \item Número de \loc{epochs}: El número de iteraciones de entrenamiento de la red. Como vimos en la teoría, hay un punto en que el \loc{accuracy} de validación de la red disminuye y comienza a sobreajustarse.

  \item \loc{Learning rate}: El learning rate nos indica lo rápido que aprende la red.
\end{itemize}

\subsection{Optimizer}

Un optimizador es un algoritmo usado en \loc{deep learning} para disminuir las funciones de pérdida. Son funciones que dependen de los parámetros de pesos, bias y en algunos casos, del \loc{learning rate}. Hay muchos optimizadores como el descenso de gradiente, \sigla{SGD} (\loc{Stochastic Gradient Descent}), \sigla{AdaGrad} (\loc{Adaptive Gradient Descent}), \sigla{RMS-Prop} (\loc{Root Mean Square Propagation}) o \sigla{Adam} (\loc{Adaptive Moment Estimation}).

En nuestro caso hemos usado \sigla{Adam} y \sigla{SGD}, siendo este último el que mejores resultados nos ha dado.

\subsection{Scheduler}

Un \loc{scheduler} se encarga de disminuir el \loc{learning rate} (\loc{lr}) en las distintas \loc{epochs} del entrenamiento.En neustro caso, el \loc{lr} utilizado es <<ReduceLROnPlateau>>. This \loc{scheduler} funciona reduciendo el \loc{lr} cuando la red deja de funcionar durante $n$ número de \loc{epochs}. El número de iteraciones que deben pasar para actualizar el \loc{lr} se denomina <<paciencia>>. En el momento en que la red no mejore el número de iteraciones establecidos, se multiplicará el valor actual del \loc{lr} por un factor de disminución.

\subsection{Criterion}

O funciones de pérdida, nos ayudan a calcular la calidad del entrenamiento de la red. Esta medida es mucho más fiable que el \loc{accuracy} global, ya que este puede estar contaminado por un valor que se está clasificando especialmente bien. A la hora de calcular la perdida entre la salida de la red y el \loc{ground truth}, buscamos que el valor sea el mínimo posible.

Las funciones de pérdida usadas son:

\begin{itemize}
  \item \loc{Cross Entropy Loss} \bib{crossentropyloss}: Esta función forma parte de las llamadas funciones de pérdida \sigla{NLL} (\loc{Negative Log-Likelihood }). Esta familia de funciones se usa únicamente en clasificación.
  \item \sigla{MSLoss} (\loc{Median Square Error Loss}) \bib{mseloss}: Esta función de pérdida se usa en la regresión de las poses. Compara el vector salido de la red con las poses correctas y calcula el error cuadrático medio.
\end{itemize}


%%%%%%%%%%%%%%%%%%%%%%%%%%%%%%%%%%%%%%%%%%%%%%%%%%%%%%%%%%%%%

\section{Dos salidas en la misma red}

En los apartados anteriores hemos hablado de clases y poses cuando nos hemos referido a la salida de la red. Esto es porque en nuestra última \sigla{CNN} creada, tenemos dos salidas, una con las probabilidades de cada clase y otra con un vector de poses sobre el gesto de entrada.

Esto significa que nos enfrentamos a dos problemas, uno de clasificación y uno de regresión. Es por esto que cada salida tiene una función de pérdida distinta.

La hipotesis que fundamento la creación de está estructura fue:

``A la hora de clasificar gestos, la red \loc{SimpleNet} se centra más en los sujetos que hacen los gestos que en los propios gestos. Se comprueba que cuando el sujeto es completamente distinto, la salida de la red identifica el gesto por la persona más que por los movimientos de las manos. Por esto se piensa que añadiendo el factor movimiento con las poses, podemos <<hacer olvidar>> a la red del sujeto y conseguir centrarla en la detección de los gestos.``

Tras pensar esto, añadimos la segunda salida de la red, creamos la fución de pérdida, y sumamos el valor del coste para dicha salida al ya existente de la clasificación. Tras esto, hacemos el \loc{backpropagation} con la suma de las pérdidas. La mejora en \loc{accuracy} es mínima pero vemos que es capaz de centrarse mucho mejor en los gestos que en la persona en concreto. Incluso es capaz de adivinar un gesto dada una entrada de un dibujo animado, como vemos en la figura \ref{fig:DrawingSign}.

\imagen{./img/memoria/aspectos/DrawingSign}{La red es capaz de adivinar un signo hecho por un dibujo animado}{DrawingSign}

Una idea muy parecida se puede ver en \bib{maruyama2021word}. En este \loc{paper} proponen usar numerosas entradas de la red junto con datos de poses. Proponen introducir la imagen completa, junto a la imagen por separado de las manos y la cara y por otro lado pretenden también introducir información de poses del cuerpo completo.

Por otro lado, en \bib{fang2017deepasl} podemos ver el acercamiento basado únicamente en poses. En este caso usa una \sigla{DRNN} (\loc{Deep Recurrent Neural Network}).

Nota: El \loc{dataset} de las poses está proporcionado por los mismo investigadores de \sigla{WLASL} \bib{li2020word}. Este dataset está a su vez generado usando \loc{PoseNet}, un modelo de generación de poses bastante renombrado \bib{cao2017realtime}\bib{wei2016cpm}.

%%%%%%%%%%%%%%%%%%%%%%%%%%%%%%%%%%%%%%%%%%%%%%%%%%%%%%%%%%%%%

\section{Exportación del modelo}

A la hora de exportar el modelo, hemos usado la herramienta interna de \prog{PyTorch}, que nos exporta modelos con formato <<\prog{.pth}>>. Este formato es perfecto si nos mantenemos en el ecosistema de \prog{PyTorch}, pero si quremos usar el modelo con otros \loc{frameworks} o librerias, debemos exportarlo a un formato estandar.

\subsection{ONNX}

El formato estandar más utilizado es \prog{ONNX}. Puedes ir a la parte de herramientas en \ref{tool:ONNX} para más información.

Con este formato, hemos podido usar el modelo con la librería de \prog{ONNXRuntime} y así conseguir una inferencia mucho más rápida. Por otro lado, hemos usado esta misma librería para comprimir el modelo usando cuantización.

%%%%%%%%%%%%%%%%%%%%%%%%%%%%%%%%%%%%%%%%%%%%%%%%%%%%%%%%%%%%%

\section{Compresión del modelo}

\subsubsection{\loc{Quantization}}

La cuantización es el proceso de reducir la precisión que tienen los pesos de las neuronas, así como los \loc{bias}, con la intención de reducir el uso de memoria y el tamaño del modelo.


Para hacer esto transformamos el tipo de dato de los pesos, que normalmente es un \loc{float} de 32 \loc{bits}, a un tipo que ocupe mucho menos espacio, como por ejemplo puede ser un \loc{int} de 8 \loc{bits}.

En este movimiento de transformar de 32 a 8 \loc{bits}, estamos reduciendo el tamaño 45 veces, reduciendo de forma significativa el consumo de memoria.

La pega que tiene esto, es que el \loc{accuracy} se puede ver ligeramente rebajado. Aquí es donde tenemos que estudiar si merece la pena la reducción de tamañoi y uso de memoria en contra de la pérdida de acierto.

En la figura \ref{fig:quantization} podemos observar la diferencia de tamaño entre el modelo en formato <<.onnx>> antes y después de cuantizar.

\imagen{./img/memoria/aspectos/quantization}{Modelo exportado antes y después de cuantizar con \prog{ONNXRuntime}}{quantization}

Un ejemplo claro de esa pérdida de \loc{accuracy} la vemos a la hora de inferir el mismo video del dibujo animado del apartado anterior. Esta vez en vez de inferir correctamente <<Book>>, infiere <<Before>>.

Por otro lado, la cuantización nos ha permitido poder subir el modelo contenedorizado de forma más sencilla y rápida.

\section{Repasemos un experimento}

Voy por último a hablar de uno de los experimentos realizados en el proyecto.

Antes de pasar a crear la \sigla{CNN} para clasificación de video, y cuando estaba en la fase en la que los datos de entrada eran únicamente \loc{frames}; pensé que podía incluir la secuencia de \loc{frames} de un video en una única imagen.

Podemos ver uno de los ejemplos en la figura \ref{fig:framesConcat}.

\imagen{./img/memoria/aspectos/framesConcat}{Concatenación de frames con la intención de representar un video (gesto es <<Before>>). Imágenes del dataset WLASL \cite{li2020word}}{framesConcat}

Esto fue un auténtico fracaso, el \loc{accuracy} de la red era mínimo, y cuando concatenabamos números \loc{frames} de una imagen, acabábamos con un tensor único enorme.

El problema principal es que una entrada de este estilo confunde mucho a una red convolucional, ya que esta intenta buscar características espaciales en la dimension, pero se está encontrando con muchos problemas:

\begin{enumerate}
  \item Píxeles <<muertos>>: Para empezar, al concatenar la imagen, hemos añadido píxeles negros entre \loc{frames}. El \loc{kernel} va a pasar por ellos y no va a recibir nin´gún tipo de información espacial.
  \item Demasiada información para una dimensión: Estamos dando información sobre las caracterísitcas espaciales de la imagen al mismo tiempo que buscamos encontrar el movimiento de los gestos entre \loc{frames}
        . Para que esto funcione correctamente, el movimiento de los gestos entre \loc{frames}
        debe estar en otra dimension, y que sea el \loc{kernel} de esa dimensión el que produzca la convolución que detecte los movimientos.
\end{enumerate}

Tras esto, pudimos comprender que necesitábamos pasarnos a la clasificación de un video, uniendo los \loc{frames}, no como una imagen única, sino como una nueva dimensión en el tensor de entrada.

\section{Estado actual del modelo}

El modelo ha sido entrenado con 8 etiquetas (\loc{all, before, book, drink, help, no, walk, yes}) y mantiene los siguientes \loc{accuracies} (tabla {tabla:accs})

\tablaSmall{Accuracy actual tra}
{c c c}{accs}
{ Entrenamiento & Test & Validación \\}{
  99,3\% & 79,88\% & 78,65\%  \\
}

Este modelo es accesible en DockerHub. Se puede consumir desde la API abierta o desde la demo en la web. El código es Open Source en Gitlab y Github.

\begin{itemize}
  \item \textbf{Dockerhub}: \url{https://hub.docker.com/repository/docker/gazquez/sign2text}
  \item \textbf{API Abierta}: \url{https://api.sign2text.com/docs}
  \item \textbf{Demo Web}: \url{https://sign2text.com/}
  \item \textbf{Repositorio del modelo}: \url{https://github.com/irg1008/Sign2Text.git}
  \item \textbf{Respositorio de la demo}: \url{https://github.com/irg1008/Sign2Text-Astro.git}
  \item \textbf{Repositorio del servidor}: \url{https://github.com/irg1008/Sign2Text-API.git}
\end{itemize}
\capitulo{6}{Trabajos relacionados}

\section{Sign2Text - \loc{Letter level translation}}

Este proyecto, con el mismo nombre que el mio, tiene 125 estrellas en Github y se trata de un traductor a tiempo real de lenguaje de signos \sigla{ASL} a texto a nivel de letra. Es capaz de traducir las letras del abecedario usando inferencia por webcam. Usa un modelo preentrenado \sigla{ResNet50} y tiene el modelo alojado en una instancia en AWS.

Podemos acceder al repositorio en \url{https://github.com/BelalC/sign2text.git}

\section{ASL Translator}

Esta aplicación, subida al App Store ya que usa el mismo \loc{dataset} que hemos usado en este proyecto, es una \loc{app} de traducción de texto a signos, es decir, todo lo contrario a la nuestra. Asegura:

\begin{enumerate}
  \item Traducir más de 30000 palabras
  \item Ser capaz de traducir frases complejas
  \item Usar un algoritmo que mezcla los vídeos con los signos de forma que muestran oraciones complejas de forma suavizada y sin cortes.
\end{enumerate}

No compite directamente con la idea de este proyecto pero parece un proyecto interesante y con un hueco en el mercado sin competencia.

Se puede acceder a la página de la \loc{app} en \url{https://apps.apple.com/us/app/asl-translator/id421784745}

\section{\loc{ASL to english} - traducción a nivel palabra}

Este proyecto, que cogió mucha fuerza después de aparecer en algún artículo y hacerse ligeramente viral en LinkedIn, tiene 674 estrellas en Github.

En este proyecto se traducen <<6>> palabras de lenguaje de signo a texto; eso sí, ninguna de las palabras se representa con un gesto en movimiento, todas son estáticas. Usa también una red preentrenada y una webcam como entrada para la inferencia.

Como curiosidad, al comienzo de este proyecto, se mantuvo una conversación con la autora por \loc{e-mail} para ver si se podría hacer una colaboración.

Se puede acceder al repositorio del proyecto en \url{https://github.com/priiyaanjaalii0611/ASL_to_English.git}


\section{How2Sign}

How2Sign \bib{Duarte_CVPR2021} es un \loc{dataset} de \sigla{ASL} no solo compuesto de imágenes y vídeos de gestos, sino que también aporta audios, transcripciones, mapas de profundidad y datos de poses.

Todos los vídeos están grabados usando un \loc{chroma key} para mayor facilidad de cambio de fondo y las poses han sido calculadas en 3D en un estudio dedicado. Por esta razón el tamaño del \loc{dataset} aumenta a 290GB.

Con una licencia \sigla{CC} y sin permiso para uso comercial, el \loc{dataset} se puede acceder en \url{https://how2sign.github.io/}
\capitulo{7}{Conclusiones y Líneas de trabajo futuras}

Vamos a analizar las conclusiones extraídas de la realización de este proyecto y ver que mejoras se pueden realizar de cara al futuro.

\section{Conclusiones}

De forma general, se han cumplido las características buscadas al inicio del proyecto.
\begin{itemize}
  \item Disposición del modelo neuronal a la comunidad. Hemos podido crear un contenedor docker con el modelo listo para ser usado. También hemos creado una \sigla{API} abierta para que cuualquier desarrollador la consuma en su aplicación de cliente. Y también hemos creado una web con una demo para que cualquier usuario lo use de forma sencilla.
  \item Se ha aprendio a usar el \loc{framework} \prog{PyTorch} de forma general, así como en detalle los módulos de \loc{supervised learning} y tratamiento y procesado de imagen y vídeo.
  \item Hemos conseguido implementar una red convolucional desde cero, comprendiendo de forma profunda su estructura y las conexiones entre las disitntas capas, así como la modificación de los vectores de salida para adaptarse a los distintos problemas.
  \item Se ha mejorado considerablemente la lectura de árticulos científicos relacionados con el campo de la ciencia de datos y el aprendizaje automático. Al finalizar, eramos capaces de analizar y comprender lo principal de un árticulo y de este modo sustraer la información útil en el desarrollo del proyecto.
  \item Se ha ampliado el número de tecnologías conocidas en tanto en el campo de la ciencia de datos como en el desarrollo \loc{software} en general.
  \item Hemos podido usar las últimas herramientas y servicio de \loc{cloud} con las que hemos posido hacer \loc{deploy} de un modelo que consume sustanciales recursos. Se ha aprendido a usar plaatformas como AWS, Azure y Google Cloud, entre otras.
\end{itemize}

Y, aunque el modelo actual basado en \sigla{CNNs} funciona  y es medianamente escalable, no es nada comparado con el \loc{State of the Art} \bib{fang2017deepasl} \bib{li2020word} en la actualidad. Por esto, veamos como podemos orientar el proyecto en el futuro.

\section{Líneas futuras}

Ha habido pasos en la creación del modelo, en los que podría haber tomado mejores decisiones. Como en la carga y transformación de los datos, en la creación de la estructura de la \sigla{CNN} o en la elección de \loc{schedulers} y \loc{optimizers}.

Por lo tanto, quedan muchas aspectos que mejorar. Algunas son los siguientes:

\begin{itemize}
  \item Cambio de \loc{dataset}: Uno de los grandes aspectos que ha ralentizado las mejoras e iteraciones de la red ha sido el \loc{dataset}. Aunque el número de vídeos del \loc{dataset} es correcto, la cantidad de vídeos por etiqueta es algo corta. Por otro lado, los datos de poses tenían muchas veces datos nulos o valores a 0, que no aportaban información en el entrenamiento. En algunas ocasiones en las que he probado el modelo con \loc{datasets} distintos, el \loc{accuracy} ha subido una media de 25\% en test y validación. Por todo esto, el primer cambio que haría sería el \loc{dataset}. el recambio perfecto es How2Sign, un \loc{dataset} mucho más extenso con datos de profundidad, gestos en 3D y transcripciones de audio \bib{Duarte_CVPR2021}.

  \item Mejorar la estructura de la red: En las redes y estructuras sobre \sigla{CNNs} ya establecidas observamos estructuras de decenas e incluso cientos de capas convolucionales, mezcladas con \loc{pooling}, seguidas de capas densas o incluso otras estructuras complejas formando un \loc{ensemble}. La idea de futuro es estudiar y completar una estructura más compleja de convolucionales, que permita escalar el modelo a un número mucho mayor de etiquetas.

  \item Probar otro tipo de redes de \loc{deep learning}, como \sigla{RNN} (\loc{Recursive Neural Netowrks})

  \item Mejorar la documentación a nivel \loc{codebase} y repositorio. Me gustaría haber podido documentar mucho mejor los resultados obtenidos. Es una pena que descubrí herramientas como \prog{Tensorboard} en las últimas etapas del proyecto.

  \item Creación de una \loc{landing page} de presentación, mejoría de la demo subida para inferencia con \loc{webcam}. Me gustaría poder mejorar el \loc{front-end} para incluir inferencia a tiempo real con el modelo en el \loc{back-end}. Es algo que puedo hacer muy pronto, con el estado actual del proyecto, pero que se ha quedado fuera de la iteración final.

  \item Utilización de \loc{criterions}, \loc{optimizers} y \loc{schedulers} distintos. No he podido probar tantas opciones como me hubiera gustado. Pretendo investigar en la literatura cuales son las que ofrecen mejor resultado en los distintos escenarios, y ver como puedo aplicarlo en este proyecto.

  \item Buscar la forma de compartirlo en la comunidad \loc{Open Source}. Me gustaría buscar la forma de expandir el proyecto, y así encontrar personas que estén interesadas en continuar con el proyecto o en colaborar de algún modo.

  \item Y mucho más. Me encantaría poder seguir mejorando el modelo e intentar terminar creando un producto viable con el que las personas se puedan beneficiar. Me parece que el objetivo es muy noble y con un poco de impulso por mi parte, o por cualquier desarrollador que siga con el proyecto; podemos conseguir algo que funcione genial.
\end{itemize}


\section{Experiencia personal}

Este proyecto ha sido una aventura y me ha abierto los ojos al mundo del \loc{deep learning} y la investigación. Cuando comencé pensaba que iba a ser capaz de crear una red neuronal con un \loc{accuracy} alto y que iba a poder crear una red \loc{transformers} con la que sintetizar las salidas de la \sigla{CNN}. También pensaba que todo esto lo iba a hacer en los cuatro primeros \loc{sprints}. Tras esto, la idea era poder crear una aplicación web que usará el modelo en los dos últimos meses.

\textbf{Nada más lejos de la realidad}.

Este proyecto ha sido una prueba a la paciencia: prueba tras prueba, cambio tras cambio. El avance era lento, y había semanas en las que me quedaba atascado cambiando la estructura de la red, o probando una y otra vez con distintos hiperparámetros.

El hecho de tratar con datos que producían un tensor de entrada de la red de 5 dimensiones ha hecho que las iteraciones y mejoras sobre el modelo sean lentas, muy lentas. En ocasiones el modelo ha estado entrenando durante casi un día completo (¡y en GPU!), inhabilitando cualquier posibilidad de trabajar en él.

Aun así, este proyecto me ha hecho comprender mucho más a fondo el comportamiento de la redes neuronales en general, y sobre todo, de las redes neuronales convolucionales. Comencé sin saber lo que era una convolución y un \loc{kernel}, con miedo a averiguar lo que significaba \loc{ResNet}; y he terminado creando una \loc{CNN} con un accuracy alto (aunque con relativas pocas etiquetas) y he perdido el miedo a investigar y leer artículos.

Por otro lado, he podido realizar una demo web para que cualquier usuario pueda probar el modelo con tan solo subir un vídeo. Esto me ha permitido aprender tecnologías web que todavía no había descubierto. También he creado una \sigla{API} en \prog{Python}, con lo que he aprendido un \loc{framework} más.

Y además me he tenido que enfrentar a los tres grandes de la computación en la nube hasta que he conseguido subir el pesado modelo a una de ellas. De esta experiencia salió una guía que se puede leer en el repositorio del servidor.

En general, la experiencia ha sido genial y el aprendizaje ha sido inmenso. Me ha hecho ver que el campo del \loc{machine learning} es el futuro en todos los ámbitos tanto en investigación como cotidianos. El uso de tecnologías que usen \loc{deep learning} está creciendo enormemente en la actualidad, y esta experiencia me impulsa a seguir investigando en este campo e intentar trabajar con y de esto en el futuro.

Tanto me ha gustado el proyecto, tanto he aprendido, y tanto me he dado cuenta que no se, que lo único que quiero hacer es seguir, y hacerlo mejor, con mejor información.



\bibliographystyle{plain}
\bibliography{bibliografia}

\end{document}
